 \documentclass[10pt]{article}
\usepackage[margin=0.8in]{geometry}
\usepackage{amsmath,amssymb,bm,mathtools}
\usepackage{booktabs,longtable,array,graphicx,float,hyperref,xcolor,enumitem}
\usepackage{siunitx}
\usepackage[skip=3pt,font=small]{caption}
\renewcommand{\topfraction}{0.92}
\renewcommand{\bottomfraction}{0.9}
\renewcommand{\textfraction}{0.06}
\renewcommand{\floatpagefraction}{0.85}
\usepackage{titlesec}
\titlespacing*{\section}{0pt}{1.2ex plus .3ex}{0.6ex}
\titlespacing*{\subsection}{0pt}{1.0ex plus .2ex}{0.4ex}
\setlist{topsep=2pt,itemsep=1pt,parsep=0pt}
\setlength{\textfloatsep}{6pt plus 2pt minus 2pt}
\setlength{\intextsep}{6pt plus 2pt minus 2pt}
\setlength{\abovedisplayskip}{4pt}
\setlength{\belowdisplayskip}{4pt}
\setlength{\abovedisplayshortskip}{2pt}
\setlength{\belowdisplayshortskip}{2pt}
\hypersetup{colorlinks=true,linkcolor=blue!50!black,citecolor=blue!50!black,urlcolor=blue!50!black}
\sisetup{group-separator={,},group-minimum-digits=4}
\newcommand{\ket}[1]{\left|#1\right\rangle}
\newcommand{\bra}[1]{\left\langle#1\right|}
\newcommand{\expect}[1]{\left\langle #1 \right\rangle}
\newcommand{\comm}[2]{\left[#1,#2\right]}
\newcommand{\argmax}{\operatorname*{arg\,max}}
\newcommand{\Var}{\operatorname{Var}}
\newcommand{\Cov}{\operatorname{Cov}}
\newcommand{\supp}{\operatorname{supp}}

\title{Part I: Measurement Cost of Generator Selection in\\Fixed-State ADAPT-VQE}
\author{[AUTHOR: names and affiliations]}
\date{September 2026}

\begin{document}
\maketitle

\begin{abstract}
This report defines and compares three measurement strategies for the generator-selection step in gradient-based ADAPT-VQE (adaptive derivative-assembled pseudo-Trotter variational quantum eigensolver): \textbf{M1: FC-UG}, fully commuting (FC) universal groups for static all-gradient estimation; \textbf{M2: BAI-FC-IG}, best-arm identification (BAI) using FC individual-gradient measurements; and \textbf{M3: BAI-FC-UG}, best-arm identification using fixed parent FC universal groups. A fourth construction, a \emph{no-sharing baseline} that uses neither mechanism, completes the two-by-two the three methods span and makes each mechanism separately attributable.

Two classes of number are reported, and they are kept apart throughout. \emph{Planning bounds} (Secs.~\ref{sec:completed}--\ref{sec:baseline}) are oracle diagnostics: exact gradients and exact fragment variances are used to compute the shots each method would need, and the elimination rule reads exact values, so it cannot misfire. A \emph{finite-shot surrogate} (Sec.~\ref{sec:finite-shot}) replaces the exact gradients with Gaussian estimates built from a finite number of shots and runs the same elimination rule on those. The surrogate samples estimators, not individual measurement outcomes, so it is not a shot-level simulation of a quantum device; the distinction and its consequences are set out in Sec.~\ref{sec:assumptions}.

\textbf{Results.} On the planning bounds, M3 is cheaper than M1 in all nine cases, by factors of $1.1$ to $192$, and cheaper than M2 in eight of nine; neither ordering is a theorem, and small problems with a nearly tied runner-up exceed M1 by up to $9\%$. Against the no-sharing baseline, elimination alone is worth $1.3$ to $253$ and sharing alone only $1.1$ to $3.4$. The planning bound resolves every M3 arm to a common radius, but on LiH and H$_2$O the shared allocation resolves the leader far more tightly, and M3 eliminating on that actual precision costs only $0.52$ to $0.74$ of its bound. On those five rows the M2/M3 ratio widens to $1.19$--$3.42$, M3 wins both stretched H$_2$O rows that the bound leaves unresolved under every grouping routine tested, and the two mechanisms compound rather than substitute (interaction $1.05$--$1.79$); on H$_4$, and on every row at bound level, they substitute. Under the finite-shot surrogate both best-arm methods cost $2.2$ to $3.0$ times their bounds on a coarse radius schedule; near the exact limit that falls to $1.03$--$1.35$ for M2 and $0.73$--$1.29$ for M3, so the inflation is discretisation, not noise, and M3 beats M1 on all nine simulated rows. Modelling the within-context correlation and eliminating on pairwise radii cuts M3's realised cost by a further $33$ to $49\%$ at no measured loss of accuracy. Every trial of the main finite-shot table returned the correct generator; the stated family-wise failure probability does not control the error rate, which stays below $2\%$ even at $\delta=0.99$.

\textbf{Scope.} Every number is a single-ADAPT-step diagnostic for one fixed state, and counts shots only: circuit depth, the number of distinct measurement settings, and state-preparation cost are excluded (Sec.~\ref{sec:assumptions}). Part III symmetry filtering and tapering are not included in the Part I tables; Part III supplies a smaller input problem to the same measurement methods.
\end{abstract}

\section{Problem statement and relation to previous work}
\label{sec:problem}

Given a molecular Hamiltonian $H$, a current ADAPT-VQE state $\ket{\psi}$, and a unitary coupled-cluster singles-and-doubles (UCCSD) generator pool $\{G_i\}$, estimate the measurement cost required to identify a generator with the largest absolute gradient
\begin{equation}
  i^\star = \argmax_i |g_i|,
  \qquad
  g_i = \bra{\psi}\comm{H}{G_i}\ket{\psi} .
  \label{eq:gradient}
\end{equation}
The scientific question is whether using winner-identification statistics and shared fully commuting measurement contexts reduces the gradient-selection measurement cost relative to static all-gradient estimation.

\textbf{The answer, in one sentence.} Mostly yes: on oracle planning bounds at family-wise $\delta=0.05$, combining both mechanisms (M3) costs $1.1$ to $192$ times less than static all-gradient estimation (M1) across the nine cases studied, and the same ordering holds on realised cost under a well-resolved schedule. The change of statistical target carries most of the saving, shared contexts being worth only $1.1$ to $3.4$ alone --- but on the larger systems sharing also makes elimination fire earlier, by resolving the leader more tightly than required, and there the two compound.

Here $g_i$ is the derivative of the energy with respect to the variational parameter of generator $G_i$ at $\theta_i=0$, in Hartree per radian; $G_i$ is anti-Hermitian, so $g_i$ is real. The generator pool and its labelling convention are fixed in Sec.~\ref{sec:computational}.

\textbf{Relation to previous work.} ADAPT-VQE and the spin-orbital UCCSD pool used here are those of Ref.~\cite{Grimsley2019ADAPT}. Two lines of work attack the selection-step measurement cost separately. Ref.~\cite{Anastasiou2023ReallyMeasure} attacks it by \emph{sharing}: it partitions the gradient observables into commuting sets using Hamiltonian terms as pivots, reducing the pool-gradient cost to a small multiple of one energy measurement, and Ref.~\cite{Shkolnikov2023Avoiding} pursues a related reduction alongside symmetry considerations. Ref.~\cite{Huang2025QuantumGambling} attacks it by \emph{elimination}: it recasts generator selection as best-arm identification and applies successive elimination, discarding unpromising candidates early. Ref.~\cite{Rossi2026RSe} changes the selection criterion itself.

The two mechanisms have not, to our knowledge, been measured against each other on one cost criterion, and it is not obvious that they compose. This report supplies the missing comparison rather than a new algorithm. Its contributions are: the $2\times2$ factorial of sharing and elimination, completed by a no-sharing baseline so that each mechanism and their \emph{interaction} are separately attributable (Sec.~\ref{sec:baseline}); a single shot criterion with no free schedule or precision parameter, applied identically to all four constructions; and the finding that, at the level of the planning bounds, the two mechanisms are substitutes rather than complements, with elimination carrying the result. M1 stands in for the sharing line, M2 for the elimination line, and M3 for their combination.

\section{Fixed method names}

The phrase ``naive/no-sharing'' is retired as a \emph{method} name. A construction that shares nothing does appear below, as the fourth corner of the two-by-two that the three methods span; it is called the \emph{no-sharing baseline} throughout, it is a diagnostic rather than a project method, and it is defined in Sec.~\ref{sec:baseline}.

\begin{description}[leftmargin=1.2em,itemsep=1pt]
  \item[\textbf{M1: FC-UG}, Fully Commuting Universal Groups.] One global parent FC grouping of the union of all Pauli products in all gradient commutators; estimate the full gradient vector.
  \item[\textbf{M2: BAI-FC-IG}, BAI with FC Individual Gradients.] Each generator gradient is a BAI arm, measured in an FC grouping of its own commutator. No cross-gradient sharing.
  \item[\textbf{M3: BAI-FC-UG}, BAI with FC Universal Groups.] The same fixed parent contexts as M1, with BAI elimination, sampling only the parent contexts needed to separate surviving generators.
  \item[\emph{No-sharing baseline}.] A diagnostic, not a project method: each generator grouped separately as in M2, every one resolved to the common radius used by M1 (Sec.~\ref{sec:baseline}).
\end{description}

\section{Object lineage}

The qubit Hamiltonian is
\begin{equation}
  H=\sum_p h_p P_p,
\end{equation}
where $P_p$ are Pauli products after Jordan--Wigner mapping. For each generator $G_i$, define the gradient operator
\begin{equation}
  C_i = \comm{H}{G_i} .
\end{equation}
The ADAPT gradient is the expectation value of this operator,
\begin{equation}
  g_i=\expect{C_i}_{\psi}.
\end{equation}
Expanding all gradient operators in one shared Pauli-coordinate system gives
\begin{equation}
  C_i=\sum_{\ell} A_{i\ell}R_{\ell},
  \qquad
  \mu_\ell=\expect{R_\ell}_{\psi},
  \qquad
  \bm g = A\bm\mu .
  \label{eq:gAmu}
\end{equation}
The parent universal Pauli support at this ADAPT step is
\begin{equation}
  \mathcal B_0=\bigcup_i \supp(C_i).
  \label{eq:B0}
\end{equation}
M1 and M3 use $\mathcal B_0$ as a shared measurement object. M2 does not; it groups each $C_i$ separately.

\section{The three algorithms}
\label{sec:algorithms}

\subsection{M1: FC-UG}

M1 estimates the full gradient vector. It forms one global FC grouping of the parent support,
\begin{equation}
  \mathcal B_0 \longrightarrow \mathcal M_1,\ldots,\mathcal M_K,
\end{equation}
where each $\mathcal M_\alpha$ is a fully commuting measurement context. Measurements of these contexts estimate the shared Pauli expectations $\mu_\ell$ and reconstruct every gradient using Eq.~\eqref{eq:gAmu}.

The implementation uses deterministic greedy FC grouping of $\mathcal B_0$ (Sec.~\ref{sec:computational}). It is a generic stand-in for the pool-wide grouping strategies of Refs.~\cite{Anastasiou2023ReallyMeasure,Shkolnikov2023Avoiding} and is not claimed to reproduce either construction; Sec.~\ref{sec:routines} shows how much the choice of routine matters.

\subsection{M2: BAI-FC-IG}

M2 changes the statistical target from full-vector estimation to winner identification. Each candidate gradient is an arm in a best-arm-identification problem. Its own commutator support is grouped separately,
\begin{equation}
  \supp(C_i) \longrightarrow \mathcal M_{i,1},\ldots,\mathcal M_{i,K_i} .
\end{equation}
Throughout, $t$ indexes the BAI round and $r$ denotes a confidence radius; the two are never the same symbol. BAI maintains confidence intervals for active gradients and eliminates a generator $i$ at round $t$ if
\begin{equation}
  |\widehat g_i|+r_i < \max_{j\in\mathcal A_t}\left(|\widehat g_j|-r_j\right),
  \label{eq:bai-elim}
\end{equation}
where $\mathcal A_t$ is the active generator set and $r_i=z_\delta\,\epsilon_i$ is the current confidence radius of arm $i$, with $\epsilon_i$ its one-standard-error target and $z_\delta$ the confidence factor of Eq.~\eqref{eq:zdelta}. All methods in this report use a single common radius $r_i=r$ across active arms; the rule is written with per-arm radii because the covariance-aware variant of Sec.~\ref{sec:correlation} does not. M2 does not exploit universal cross-gradient sharing.

\textbf{Implementation status.} M2 is now implemented as its own module. Each run writes the individual-gradient FC grouping $\mathcal M_{i,k}$ for every generator, the oracle fragment standard deviation of every fragment, the per-arm shot cost, and the full BAI elimination history with confidence radii and active sets. The M2 column is therefore auditable; it remains an oracle planning diagnostic rather than a finite-shot result.

\subsection{M3: BAI-FC-UG}

M3 combines the mechanisms of M1 and M2. It starts with the same fixed parent universal FC contexts as M1,
\begin{equation}
  \mathcal M_1,\ldots,\mathcal M_K,
\end{equation}
but BAI eliminates generator candidates,
\begin{equation}
  \mathcal A_0 \supset \mathcal A_1 \supset \mathcal A_2 \supset \cdots .
\end{equation}
At BAI round $t$, the active Pauli support is
\begin{equation}
  \mathcal B_t=\bigcup_{i\in\mathcal A_t}\supp(C_i),
  \qquad
  \mathcal B_t\subseteq\mathcal B_0 .
\end{equation}
The active measurement contexts are restricted parent contexts,
\begin{equation}
  \mathcal M_\alpha^{(t)}=\mathcal M_\alpha\cap\mathcal B_t .
\end{equation}
The parent contexts are not rebuilt during BAI. This nested construction keeps previous measurements compatible while allowing the algorithm to stop sampling information that only affects eliminated generators.

\textbf{How M3's cost is aggregated.} Shots taken in a context are reusable, so a context is charged once, at the largest number of shots any round ever demanded of it. Writing $n_\alpha^{(t)}$ for the allocation of Eq.~\eqref{eq:allocation} solved over the active set $\mathcal A_t$ at the radius of round $t$, the reported M3 cost is
\begin{equation}
  M_{\mathrm{M3}}=\sum_{\alpha}\max_{t} n_\alpha^{(t)} .
  \label{eq:m3-aggregate}
\end{equation}
Each round's allocation is solved independently, without reference to shots already committed; the maximum is taken afterwards. This is the accounting a real run would get for free by keeping its data, and it is why the trajectory is fully determined by the finitely many rounds at which the active set changes.

\section{Oracle shot proxy}
\label{sec:proxy}

All numbers in Secs.~\ref{sec:completed}--\ref{sec:baseline} are oracle planning estimates: the planner is given the exact $g_i$ and $\sigma_{i\alpha}$, and nothing is learned from shots. \emph{Oracle} and \emph{noiseless} mean only that. They do not refer to a noiseless device, and state-preparation error is excluded from the model altogether (Sec.~\ref{sec:assumptions}).

\textbf{Sorted order and the gap.} Order the pool by decreasing absolute gradient, $|g_{(1)}|\ge|g_{(2)}|\ge\cdots$; the parenthesised index is a rank, never a generator index. Then
\begin{equation}
  \mathrm{gap}=|g_{(1)}|-|g_{(2)}|,
  \qquad
  \Delta_i=|g_{(1)}|-|g_i| \ \ge 0 ,
  \label{eq:gap}
\end{equation}
so $\Delta_{(2)}=\mathrm{gap}$ and $\Delta_{(1)}=0$.

\textbf{Shots and context-shots.} One \emph{shot} is one state preparation plus one measurement of one context: a single circuit execution. \emph{Context-shot} is the same thing charged to the context it was spent in; every cost quoted is $\sum_\alpha n_\alpha$. The metric counts shots only --- not circuit depth, not the number of \emph{distinct} measurement settings, not classical post-processing (Sec.~\ref{sec:assumptions}).

For one gradient operator written as a sum of independently measured FC fragments,
\begin{equation}
  C_i = \sum_{\alpha=1}^{K_i} F_{i\alpha},
\end{equation}
with oracle fragment variances
\begin{equation}
  \sigma^2_{i\alpha}=\Var_\psi(F_{i\alpha})
  =\expect{F_{i\alpha}^2}_\psi-\expect{F_{i\alpha}}_\psi^2,
\end{equation}
the optimal independent-fragment allocation gives the approximate total number of shots needed to estimate $g_i$ with one-standard-error target $\epsilon_i$ as
\begin{equation}
  M_i(\epsilon_i)\approx
  \frac{\left(\sum_{\alpha=1}^{K_i}\sigma_{i\alpha}\right)^2}{\epsilon_i^2} .
  \label{eq:single-gradient-shot}
\end{equation}
\textbf{Precision target and confidence radius.} $\epsilon_i$ is a one-standard-error target on $\widehat g_i$; a \emph{confidence radius} is the half-width of the corresponding interval,
\begin{equation}
  r_i=z_\delta\,\epsilon_i ,
  \label{eq:radius}
\end{equation}
so a method asked for radius $r$ must reach $\epsilon=r/z_\delta$, and $\epsilon_i$ in Eq.~\eqref{eq:single-gradient-shot} is replaced by $r_i/z_\delta$. The confidence factor is the two-sided normal quantile with a Bonferroni correction over the whole pool,
\begin{equation}
  z_\delta=\Phi^{-1}\!\left(1-\frac{\delta}{2N}\right),
  \qquad N=|\{G_i\}| ,
  \label{eq:zdelta}
\end{equation}
with $N$ the \emph{full} pool size --- $26$, $92$, $140$ --- not the non-zero count and not the active set, and no correction over rounds. At $\delta=0.05$ this is $z_\delta=3.10$, $3.46$ and $3.57$, applied identically to all four constructions so that it cancels from every planning-bound ratio (Sec.~\ref{sec:sensitivity}). $\delta=0.05$ is the conventional level; Bonferroni is chosen as the simplest family-wise correction valid under arbitrary dependence between the arms, which matters here because M1 and M3 read several arms off the same shots; Holm's step-down procedure is equally valid and uniformly more powerful, and would be the better default. The price is conservatism, and Sec.~\ref{sec:calibration} measures it: the construction does not control the realised selection-error rate, so it is a convention rather than a guarantee.

Eq.~\eqref{eq:single-gradient-shot} is applied separately inside each gradient arm, because M2 measures each arm in its own contexts. M1 and M3 instead share contexts between generators, so the allocation over contexts is not separable and is obtained by solving the convex programme
\begin{equation}
  \min_{n_\alpha \ge 0} \ \sum_{\alpha} n_\alpha
  \qquad\text{subject to}\qquad
  \Var(\widehat g_i)=\sum_{\alpha}\frac{\sigma_{i\alpha}^{2}}{n_\alpha}\le\epsilon^{2}
  \quad\text{for every active } i ,
  \label{eq:allocation}
\end{equation}
where $n_\alpha$ is the number of shots taken in context $\alpha$ and $\Var(\widehat g_i)=\sum_\alpha\sigma_{i\alpha}^2/n_\alpha$ is the \emph{reconstruction variance} of gradient $i$ under that allocation: the variance of the estimate obtained by reading $\widehat g_i$ off the shared contexts through Eq.~\eqref{eq:gAmu}. The quantity $\sum_\alpha\sigma_{i\alpha}^2$ that appears repeatedly below is its value at one shot per context, and serves as a scale-free measure of how expensive generator $i$ is to reconstruct from the parent grouping. Its stationarity condition is $n_\alpha=\sqrt{\sum_i\lambda_i\sigma_{i\alpha}^{2}}$ for multipliers $\lambda_i\ge0$, solved by a damped multiplicative fixed-point iteration (damping $0.5$, at most $500$ iterations, stopping when the least-loaded constraint is within $10^{-9}$ of the tightest) and rescaled so the tightest constraint holds exactly. The rescaling makes the result feasible whatever the iteration converged to, so the stopping rule affects optimality and never validity. It does affect the last digits: rebuilt in a second numerical environment, M1 and M2 reproduce exactly on every row, while M3, whose total collects many allocations, moves by $0.1$ to $0.6\%$ on four rows (always downward). Table~\ref{tab:completed} quotes the original run; sections that recompute M3 quote the recomputation, and no ratio moves by more than $0.02$. For a single gradient it reduces exactly to Eq.~\eqref{eq:single-gradient-shot}, and it matches a general convex solver on the instances tested (Sec.~\ref{sec:tests}).

\textbf{Shot counts are not rounded.} $n_\alpha$ is real-valued with no per-context minimum; the per-arm closed form of Eq.~\eqref{eq:m2-exact} is rounded up per arm. A floor of one shot per context would add at most $K$ shots, at most $0.2\%$ of any entry in Table~\ref{tab:completed} (worst case: $53$ contexts under M3's $33{,}432$ on H$_4$ side $2.0$ CISD), changing quoted digits but no ratio.

\textbf{The schedule-free limit.} As the common radius $r$ decreases continuously, generator $i$ leaves the active set the moment $2r$ falls below $\Delta_i$, so it is resolved to exactly $\Delta_i/2$ and the leader survives to $\mathrm{gap}/2$. Measurements are reusable, so an arm leaving at radius $r$ has paid for radius $r$ and no more, and a shared context has paid the largest requirement it ever saw. The active set changes only at the distinct values of $\Delta_i$, so evaluating the allocation there determines the whole trajectory; for M2 it collapses to Eq.~\eqref{eq:m2-exact}. This removes the free shrink factor (Sec.~\ref{sec:sensitivity}).

Because these formulae depend on the exact grouping and fragment variances, every production run saves the FC groups, the fragment variances and the BAI histories used to generate each table entry.

\section{Computational details}
\label{sec:computational}

\subsection{Molecular input}

All cases use the minimal STO-3G basis, charge $0$, and a closed-shell singlet restricted Hartree--Fock reference, with no frozen core and no qubit tapering. Water is placed in the $xz$ plane with the oxygen at the origin and the bond angle held at $104.52^\circ$ while $R(\mathrm{OH})$ is varied, so the stretch is symmetric. The set is chosen to vary pool size at fixed basis and method --- $26$, $92$ and $140$ generators on $8$, $12$ and $14$ qubits --- while staying small enough that the exact FCI energy and exact fragment variances are available, since every number here depends on having them. The stretched geometries (LiH at $3.0$~\si{\angstrom}, H$_4$ at side $2.0$~\si{\angstrom}, H$_2$O at $1.75$~\si{\angstrom}) put each system in a strongly correlated regime. On H$_4$ and H$_2$O this also crowds the leading gradients together and makes selection hard; on LiH it does not, its gap of $0.096$ being among the widest here, and with no LiH equilibrium row there is nothing to compare it against. The two H$_2$O geometries come closest to isolating the effect at fixed pool size, but not cleanly: the orbital reordering documented in Sec.~\ref{sec:structure} means their Pauli-level structures differ. Four systems cannot separate size from geometry and state, and no claim in this report about a trend with size should be read as established.

Geometries, in \si{\angstrom}: LiH is Li at the origin and H at $(0,0,3.0)$,
$6$ spatial orbitals. H$_4$ is a planar square in the $xy$ plane, H at $(0,0,0)$,
$(s,0,0)$, $(0,s,0)$, $(s,s,0)$ with $s=1.0$ and $2.0$, $4$ orbitals. H$_2$O has
O at the origin and H at $(\pm R\sin\tfrac{\theta}{2},0,R\cos\tfrac{\theta}{2})$
with $\theta=104.52^\circ$ and $R=0.9572$ and $1.75$, $7$ orbitals.

Two states are used per geometry. \textbf{HF} is the closed-shell Hartree--Fock determinant. \textbf{CISD} (configuration interaction with singles and doubles) is the lowest eigenvector of the qubit Hamiltonian restricted to the determinants within double excitation of that reference, in the target $(N_\alpha,N_\beta)$ sector, normalised to unity and embedded in the full $2^{n}$-dimensional space. The CISD state is used as a correlated proxy for a later ADAPT state: it is a correlated wavefunction in the same sector, cheap to construct exactly, and it shifts the gradient spectrum in the direction a converging ADAPT run does --- the top gradient falls by a factor of $2.1$ to $44$ (the $|g_{(1)}|$ values are listed with Table~\ref{tab:completed}), which is what makes selection expensive late in a trajectory. That is the whole of the justification. It is not a saved trajectory snapshot, and no evidence is offered that its gradient \emph{spectrum} resembles one; testing that needs trajectory data (Sec.~\ref{sec:remaining}). LiH is run with the HF reference only; there is no reason of principle for the omission.

\subsection{Generator pool, conventions and labels}

Spin orbitals are interleaved: spatial orbital $p$ gives spin orbitals $2p$ ($\alpha$) and $2p+1$ ($\beta$), zero-based. The computational basis is indexed big-endian, matching the sparse-operator convention of the linear-algebra backend. The Jordan--Wigner transform is used throughout.

The generator universe is the standard spin-orbital UCCSD pool. Each generator is anti-Hermitian,
\begin{equation}
  G_i=\mathcal{JW}\!\left(\tau_i-\tau_i^\dagger\right),
  \qquad
  \tau_{i}^{\text{S}}=a_a^\dagger a_i,
  \qquad
  \tau_{i}^{\text{D}}=a_a^\dagger a_b^\dagger a_j a_i ,
\end{equation}
so that $g_i=\expect{\comm{H}{G_i}}_\psi$ of Eq.~\eqref{eq:gradient} is real and equals $\partial E/\partial\theta_i$ at $\theta_i=0$ for the parameterisation $e^{\theta_i G_i}$. Singles run over occupied--virtual pairs of the same spin; doubles conserve $N_\alpha$ and $N_\beta$ separately. The pool therefore conserves electron number and $S_z$ but not generally $S^2$. Ordering is deterministic: all singles in ascending $(i,a)$, then all doubles in ascending $(i,j,a,b)$. The labels used in the tables are exactly this: \texttt{S i -> a} for a single and \texttt{D i,j -> a,b} for a double, with zero-based interleaved spin-orbital indices. The resulting pool sizes are $26$, $92$ and $140$ for H$_4$, LiH and H$_2$O. This is the pool of the original ADAPT-VQE construction~\cite{Grimsley2019ADAPT}, chosen so that the observable structure in Sec.~\ref{sec:structure} is the standard one; qubit and spin-adapted pools give different supports and different pool sizes, and the comparison has not been repeated for them.

\subsection{Fully commuting grouping routine}

FC group counts are routine-dependent, and two rows are comparable only if produced by the same deterministic routine with the same Pauli ordering, tie-breaking rule and seed. Every row in the main tables uses one routine: deterministic greedy first-fit grouping, with Pauli strings sorted by decreasing weight (the number of non-identity factors) and ties broken by ascending lexicographic order of the string read left to right from qubit $0$, under the letter order $\mathrm{I}<\mathrm{X}<\mathrm{Y}<\mathrm{Z}$, with each candidate joining the lowest-numbered existing group that accepts it and opening a new group otherwise. No random seed is used. Sec.~\ref{sec:routines} varies this convention deliberately and reports the spread.

Fully commuting grouping is chosen over qubit-wise commuting because it admits at least as many terms per context, so it tends to reduce the number of contexts and with it the quantity this report counts --- though not necessarily, since positive covariance between terms inside a context raises that context's variance. The cost of that choice is a Clifford diagonalisation circuit per context, which the shot metric does not price (Sec.~\ref{sec:assumptions}); a qubit-wise-commuting comparison would trade more contexts for cheaper circuits and is the natural companion study once circuits are priced.

\subsection{Validation gate}

Every case is checked before any gradient is computed: the Jordan--Wigner (JW) Hamiltonian must reproduce the PySCF restricted Hartree--Fock (RHF) energy on the Hartree--Fock (HF) determinant and the PySCF full configuration interaction (FCI) energy as its lowest eigenvalue in the $(N_\alpha,N_\beta)$ sector, both to $10^{-8}$ Ha. All nine cases pass, with a maximum observed error of $4\times10^{-15}$ Ha. A case that fails the gate aborts rather than producing numbers.

\subsection{Non-zero gradient threshold}

Gradients that vanish by symmetry appear at the $10^{-8}$ numerical noise floor of the integral and SCF pipeline, while the smallest physically non-zero gradient in these cases is of order $10^{-3}$. The counts below use a threshold of $10^{-6}$, which lies on the plateau between the two. The threshold is not innocuous at tighter values: for LiH the count is 26 at $10^{-6}$ and 34 at $10^{-9}$, and every run records the count at $10^{-4}$, $10^{-6}$, $10^{-8}$ and $10^{-9}$ so the choice stays visible. \textbf{The threshold is used for reporting only.} No method filters the pool by it: all four constructions carry every one of the $N$ generators as an arm, and $N$ in Eq.~\eqref{eq:zdelta} is the full pool size. The one place the distinction is used analytically is the diagnostic in Sec.~\ref{sec:m1-cost}, which asks what M1 would cost if it were given the non-zero set for free.

\section{What the model assumes and excludes}
\label{sec:assumptions}

\textbf{Statistical assumptions, in one place.} (i) Estimators are Gaussian; the normal limit is assumed, not tested against shot-level sampling of measurement outcomes. (ii) Distinct measurement contexts are sampled independently. (iii) Within a context, the planning bounds and the main finite-shot tables treat different generators' fragments as independent; this is exact for M2 and the baseline, which give every generator its own shots, and an approximation for M1 and M3. Sec.~\ref{sec:correlation} removes it and reports what changes. (iv) State preparation is exact and noiseless, and there is no gate, readout or decoherence error anywhere in the model. (v) The union bound behind Eq.~\eqref{eq:zdelta} is taken over the $N$ pool members only; it is not taken over the sequence of elimination rounds, which in the exact limit is a continuum. (vi) Fragment variances are the exact $\Var_\psi(F_{i\alpha})$, so no cost is charged for learning them. (vii) In the finite-shot surrogate the fluctuation of each context is carried as a running sum over the shots taken in it, so successive rounds are correlated as real accumulated data would be; what is not modelled is the cost of learning the variances, per (vi).

\textbf{``System size'' means the pool size} $N$, ordering H$_4$ $(26)<$ LiH
$(92)<$ H$_2$O $(140)$. On these systems that agrees with qubit count
$(8,12,14)$ and $|\mathcal B_0|$ $(1{,}276;\ 30{,}196;\ 81{,}566)$ but not with
electron count. With four systems, and geometry and state varying with them, no
claim here about a trend with size is established; the phrase labels the
ordering.

\textbf{What the cost metric leaves out.} Circuit depth and two-qubit count for
the diagonalising Cliffords, compilation and calibration of distinct measurement
settings, and classical post-processing. The number of distinct settings differs
sharply between the methods, against the individually grouped ones.

\begin{table}[!htbp]
\centering
\scriptsize
\caption{Distinct measurement settings each construction requires, a cost component the shot metric does not price. M1 and M3 share one parent grouping; M2 and the baseline group every generator separately, so their settings are the sum over arms and are not reused between arms.}
\label{tab:settings}
\begin{tabular}{lrrr}
\toprule
Case & M1, M3: $K$ & M2, baseline: $\sum_i K_i$ & Largest single arm: $\max_i K_i$ \\
\midrule
LiH, $R=3.0$~\si{\angstrom} & 980 & 5,744 & 101 \\
H$_4$ square, side $1.0$ and $2.0$~\si{\angstrom} & 53 & 183 & 10 \\
H$_2$O eq., $R(\mathrm{OH})=0.9572$~\si{\angstrom} & 2,366 & 13,235 & 159 \\
H$_2$O stretch, $R(\mathrm{OH})=1.75$~\si{\angstrom} & 2,356 & 14,263 & 162 \\
\bottomrule
\end{tabular}
\end{table}

M2 and the baseline need $3.5$ to $6.1$ times as many distinct settings as M1 and M3 here. Nothing in this report prices that, and a metric that did could move the M2-versus-M3 comparison, which on shots alone is already within a factor of two on most rows. The counts are structural properties of the groupings, not measured circuit costs.

\section{Observable structure for all systems}
\label{sec:structure}

\begin{table}[!htbp]
\centering
\scriptsize
\caption{Global universal support used by M1: FC-UG and M3: BAI-FC-UG. The FC contexts are global parent contexts for $\mathcal B_0$, not per-gradient groupings, and all rows use the single greedy routine fixed in Sec.~\ref{sec:computational}. Non-zero gradient counts use the $10^{-6}$ threshold. Distances in \si{\angstrom}.}
\label{tab:structure}
\begin{tabular}{llrrrrr}
\toprule
Case & State & Qubits & Pool & Universal terms & FC-UG contexts & Nonzero $g_i$ \\
\midrule
LiH, $R{=}3.0$ & HF & 12 & 92 & 30,196 & 980 & 26 \\
H$_4$, side $1.0$ & HF & 8 & 26 & 1,276 & 53 & 10 \\
H$_4$, side $1.0$ & CISD & 8 & 26 & 1,276 & 53 & 10 \\
H$_4$, side $2.0$ & HF & 8 & 26 & 1,276 & 53 & 10 \\
H$_4$, side $2.0$ & CISD & 8 & 26 & 1,276 & 53 & 10 \\
H$_2$O eq., $R{=}0.9572$ & HF & 14 & 140 & 81,566 & 2,366 & 40 \\
H$_2$O eq., $R{=}0.9572$ & CISD & 14 & 140 & 81,566 & 2,366 & 48 \\
H$_2$O str., $R{=}1.75$ & HF & 14 & 140 & 81,566 & 2,356 & 40 \\
H$_2$O str., $R{=}1.75$ & CISD & 14 & 140 & 81,566 & 2,356 & 48 \\
\bottomrule
\end{tabular}
\end{table}

All FC context counts come from the single routine fixed in Sec.~\ref{sec:computational}.

\textbf{Why the two H$_2$O geometries have the same support size but different context counts.} The grouping routine is deterministic in the set of Pauli strings alone, so equal supports would have to give equal counts; $2{,}366$ against $2{,}356$ therefore means the supports are not equal. They are not. Both contain $81{,}566$ products, and $\sum_i|\supp(C_i)|$ is $314{,}720$ in both, but their intersection is only $78{,}113$: each geometry has $3{,}453$ Pauli products the other lacks, a symmetric difference of $6{,}906$, or $8.5\%$ of each. The differing products are not numerical dust --- their largest commutator coefficients reach $0.09$ and $0.11$, and their medians, $2.6\times10^{-3}$ and $3.7\times10^{-3}$, are within a factor of two of the median over the shared products.

The cause is an orbital reordering. Both geometries have the same point group and
the same multiset of orbital symmetries, but the orbital energies cross: the
canonical order is $(A_1,A_1,B_2,A_1,B_1,A_1,B_2)$ at $0.9572$~\si{\angstrom} and
$(A_1,A_1,B_1,B_2,A_1,A_1,B_2)$ at $1.75$. The symmetry-allowed integrals are
equal in number, which is why $|\mathcal B_0|$ and $\sum_i|\supp(C_i)|$ agree
exactly, but they sit at different orbital indices and so map to different Pauli
strings. The same reordering moves the top generator from \texttt{D 4,5 $\to$
12,13} to \texttt{D 6,7 $\to$ 12,13}. The two geometries are comparable as
physics, their Pauli-level structures are not identical, and a difference of ten
contexts carries no information.

\section{The complete comparison}
\label{sec:completed}

All four constructions are compared at one criterion: the winner must be
separated from the rest of the pool at family-wise confidence $\delta=0.05$. M1
and the baseline reach that with a single common radius $\mathrm{gap}/2$ over the
whole pool; M2 and M3 reach it by elimination, resolving each generator only to
$\Delta_i/2$. No free schedule or precision parameter remains in any column. The
cost scales as the inverse square of the relevant deficit, so the gap column is
the main driver of difficulty.

\textbf{M1 is given the exact gap, and this favours it.} A non-adaptive method
cannot know $\mathrm{gap}$ before measuring anything, yet the target radius
$\mathrm{gap}/2$ is computed from the exact gradients. M1 is therefore charged
the loosest radius that happens to suffice, the most generous reading available.
On H$_4$ side $1.0$ HF the M1/M2 ratio moves from $0.66$ to $168$ as the target
radius $\mathrm{gap}/d$ varies over $d=1,2,4,8,16$, M1 costing $561{,}216$,
$2{,}244{,}862$, $8{,}979{,}447$, $35{,}917{,}787$ and $143{,}671{,}148$ shots. The same oracle knowledge
sets the good-enough tolerance of Sec.~\ref{sec:good-enough}. Read every M1 column
as \emph{M1 at its best}.

\begin{table}[!htbp]
\centering
\scriptsize\setlength{\tabcolsep}{3.2pt}
\caption{Oracle context-shot comparison at the winner-identification criterion, $\delta=0.05$. M1/M2 $>1$ means M2 is cheaper than M1; M2/M3 $>1$ means M3 is cheaper than M2. Entries below unity are bold. Distances in \si{\angstrom}; the baseline is the no-sharing diagnostic of Sec.~\ref{sec:baseline}. $^\ast$Gap to four significant figures; full values follow the table.}
\label{tab:completed}
\begin{tabular}{llrrrrrrrr}
\toprule
Case & State & Gap$^\ast$ & Baseline & M1 & M2 & M3 & M1/M2 & M2/M3 & M1/M3 \\
\midrule
LiH, $R{=}3.0$ & HF & 0.0964 & 6,819,980 & 3,563,904 & 1,140,583 & 638,423 & 3.12 & 1.79 & 5.58 \\
H$_4$, side $1.0$ & HF & 0.0215 & 4,848,846 & 2,244,862 & 854,629 & 566,951 & 2.63 & 1.51 & 3.96 \\
H$_4$, side $1.0$ & CISD & 0.1181 & 130,375 & 59,642 & 102,806 & 53,727 & \textbf{0.58} & 1.91 & 1.11 \\
H$_4$, side $2.0$ & HF & 0.0166 & 2,002,941 & 598,102 & 422,746 & 220,389 & 1.41 & 1.92 & 2.71 \\
H$_4$, side $2.0$ & CISD & 0.0519 & 180,604 & 58,068 & 86,857 & 33,432 & \textbf{0.67} & 2.60 & 1.74 \\
H$_2$O eq., $R{=}0.9572$ & HF & 0.1011 & 204,022,319 & 172,147,028 & 25,096,606 & 22,130,852 & 6.86 & 1.13 & 7.78 \\
H$_2$O eq., $R{=}0.9572$ & CISD & 0.00206 & 485,001,056,522 & 411,524,732,152 & 46,457,815,707 & 44,971,732,317 & 8.86 & 1.03 & 9.15 \\
H$_2$O str., $R{=}1.75$ & HF & 0.0219 & 4,407,619,149 & 3,864,803,902 & 17,423,631 & 20,163,542 & 221.81 & \textbf{0.86} & 191.67 \\
H$_2$O str., $R{=}1.75$ & CISD & 0.0170 & 6,317,143,181 & 5,715,348,735 & 168,305,383 & 165,843,674 & 33.96 & 1.01 & 34.46 \\
\bottomrule
\end{tabular}
\end{table}

The winners are \texttt{D 2,3 $\to$ 10,11} (LiH), \texttt{D 2,3 $\to$ 6,7} (both
H$_4$ HF rows), \texttt{D 0,1 $\to$ 6,7} (both H$_4$ CISD rows), \texttt{D 4,5
$\to$ 12,13} (H$_2$O equilibrium) and \texttt{D 6,7 $\to$ 12,13} (H$_2$O
stretched), with $|g_{(1)}|$ of $0.243187$, $0.295316$, $0.142467$, $0.315036$,
$0.098961$, $0.304910$, $0.006925$, $0.388493$ and $0.106873$ down the rows. All
four constructions return the true maximiser in all nine cases, which in oracle
mode is a consistency check rather than a result: the rule reads exact gradients,
so it cannot eliminate the maximiser. Whether it is still found once the
gradients are estimated is the subject of Sec.~\ref{sec:finite-shot}.

M3 is cheaper than M1 in every row. M3 cheaper than M2 in eight rows of nine and
M2 than M1 in seven, both subject to the routine-dependence of
Sec.~\ref{sec:routines}.

\subsection{M3 is not guaranteed to beat M1}
\label{sec:m3-m1}

The intuition for $M_{\mathrm{M3}}\le M_{\mathrm{M1}}$ is that M3 solves the same
allocation over a shrinking subset of the constraints at radii never tighter than
M1's, so every breakpoint's allocation costs less than M1's in total. But the M3
cost is the sum over contexts of the per-context \emph{maximum} across breakpoints,
Eq.~\eqref{eq:m3-aggregate}, and two breakpoints can load different contexts.
\textbf{The ordering is not a theorem.} An adversarial search over small random
problems --- sparse fragment standard deviations and random gradients, hill-climbed
on $M_{\mathrm{M3}}/M_{\mathrm{M1}}$ --- ends above unity from $19$ of $60$
unbiased starts and $24$ of $60$ starts biased towards a crowded top of the
spectrum, reaching $1.09$ in both. Every allocation in the worst instance agrees
with a brute-force optimum to $10^{-4}$, so the excess is real, not a solver
artefact.

The mechanism is specific. In the worst instance, four arms over two contexts, the
two runners-up are nearly tied, so the last two breakpoints differ in radius by
$0.2\%$; but the arm that leaves between them was pinning the allocation to
context A, and without it the leader's optimum moves about seventy shots to
context B. Taking the maximum per context pays for both loadings. The excess is
therefore a property of the accounting --- each breakpoint solved independently of
the shots already committed --- rather than of the method: topping up from the
shots already taken instead gives $216.3$ against M1's $216.8$ on the same
instance. None of the nine cases comes close; the tightest is H$_4$ side $1.0$
CISD at $M_{\mathrm{M1}}/M_{\mathrm{M3}}=1.12$, whose top of the spectrum is also
the most crowded in the set.

\subsection{What sets the M1 cost}
\label{sec:m1-cost}

M1 must reach its common radius for every generator simultaneously, so its cost is set by the pool member with the largest reconstruction variance $\sum_\alpha\sigma_{i\alpha}^2$ --- the \emph{binding generator}, in the sense that its constraint in Eq.~\eqref{eq:allocation} is the one the final rescaling makes tight. Table~\ref{tab:binding} identifies it in every case.

\begin{table}[!htbp]
\centering
\scriptsize
\caption{The generator whose constraint binds M1, against the winner. $\sum_\alpha\sigma^2$ is the reconstruction variance at one shot per parent context. The last column is the factor by which M1's cost falls if it is handed the set of non-zero-gradient generators for free --- information no non-adaptive method has.}
\label{tab:binding}
\begin{tabular}{llrrrr}
\toprule
Case & Binding generator & $|g|$ of it & $\sum_\alpha\sigma^2$ & same, winner & M1 all / M1 non-zero \\
\midrule
LiH, $R{=}3.0$, HF        & \texttt{D 0,1 $\to$ 4,9} & $0$                 & 11.76  & 0.072 & 2.28 \\
H$_4$, side $1.0$, HF     & \texttt{D 0,1 $\to$ 4,7} & $0$                 & 2.713  & 0.447 & 1.71 \\
H$_4$, side $1.0$, CISD   & \texttt{D 0,1 $\to$ 6,7} & $1.4\times10^{-1}$  & 1.741  & 1.741 & 1.56 \\
H$_4$, side $2.0$, HF     & \texttt{D 0,1 $\to$ 4,7} & $0$                 & 0.218  & 0.067 & 1.83 \\
H$_4$, side $2.0$, CISD   & \texttt{S 1 $\to$ 7}     & $1.9\times10^{-31}$ & 0.144  & 0.122 & 1.93 \\
H$_2$O eq., HF            & \texttt{S 0 $\to$ 10}    & $7.1\times10^{-14}$ & 537.4  & 1.317 & 2.20 \\
H$_2$O eq., CISD          & \texttt{S 0 $\to$ 10}    & $1.6\times10^{-5}$  & 530.8  & 1.284 & 2.01 \\
H$_2$O str., HF           & \texttt{S 0 $\to$ 10}    & $1.8\times10^{-8}$  & 531.6  & 0.592 & 2.35 \\
H$_2$O str., CISD         & \texttt{S 0 $\to$ 10}    & $2.0\times10^{-4}$  & 455.9  & 0.565 & 2.02 \\
\bottomrule
\end{tabular}
\end{table}

In six of the nine cases the binding generator has a gradient below the $10^{-6}$ threshold, and its reconstruction variance dwarfs the winner's --- by a factor of 164 on LiH and of 400 to 900 on H$_2$O. In the remaining three it carries a small but non-zero gradient, and on H$_4$ side $1.0$ with the CISD state it is the winner itself.

This is not a defect of the comparison: M1 is non-adaptive, so it must bound every generator tightly enough to rule it out, and paying for high-variance operators that turn out not to matter is precisely the inefficiency elimination removes. Nor is it the whole story. Restricting M1 to the non-zero-gradient generators --- the last column of Table~\ref{tab:binding} --- cuts its cost by only $1.6$ to $2.4$, because the rest also carry substantial variance. The cost is spread over the pool rather than concentrated in one irrelevant operator. This is also a rough guide to how much Part III symmetry filtering could do for M1, although not a bound: filtering and tapering transform the operators as well as dropping generators, so the two are not the same operation.

\subsection{Where the cost actually sits: breadth and depth}
\label{sec:breadth-depth}

The shot cost of a BAI method is not concentrated in the final elimination, as a
naive reading of the $1/r^2$ scaling suggests. It splits into two terms, and
which one dominates changes with the system.

\textbf{One definition, applied the same way to both methods.} Let
$r_0=\tfrac12\max_i\Delta_i$ be the coarsest radius of the trajectory, the radius
at which the first generator is eliminated and at which every generator is still
active. \emph{Breadth} is the cost of that first sweep,
\begin{equation}
  M^{\text{breadth}} = \text{cost of resolving every generator to } r_0 ,
  \qquad
  M^{\text{depth}} = M^{\text{total}} - M^{\text{breadth}} ,
  \label{eq:breadth-depth}
\end{equation}
and the \emph{breadth share} is $M^{\text{breadth}}/M^{\text{total}}$. For M2 the
cost is separable, so $M^{\text{breadth}}=\sum_i (z_\delta S_i^{\text{own}}/r_0)^2$,
writing $S_i^{\text{own}}=\sum_{\alpha}\sigma_{i\alpha}$ over generator $i$'s own
contexts and $S_i^{\text{par}}$ for the same sum over the parent contexts. For M3 it is
the total of the allocation at the first breakpoint, before any elimination.
Because shots are reusable, breadth is a cost neither method can avoid: it is the
price of learning which generators are worth pursuing.

\begin{table}[!htbp]
\centering
\scriptsize
\caption{Breadth share, Eq.~\eqref{eq:breadth-depth}: the fraction of each method's total that is already committed at the coarsest radius, when every generator is still active. The remainder is the depth term. Both columns use the same definition; M3's is the cumulative allocation at the first breakpoint of its elimination history.}
\label{tab:breadth-depth}
\begin{tabular}{lrrr}
\toprule
Case & M2 breadth share & M3 breadth share & $\max_i\Delta_i/\mathrm{gap}$ \\
\midrule
H$_4$, side $1.0$, HF   & 3.0\%  & 2.1\%  & 13.7 \\
H$_4$, side $1.0$, CISD & 87.1\% & 76.3\% & 1.2 \\
H$_4$, side $2.0$, HF   & 1.3\%  & 0.8\%  & 18.9 \\
H$_4$, side $2.0$, CISD & 57.1\% & 47.7\% & 1.9 \\
LiH, $R{=}3.0$, HF      & 93.9\% & 87.7\% & 2.5 \\
H$_2$O eq., HF          & 89.4\% & 85.5\% & 3.0 \\
H$_2$O eq., CISD        & 92.8\% & 81.3\% & 3.4 \\
H$_2$O str., HF         & 80.1\% & 60.7\% & 17.8 \\
H$_2$O str., CISD       & 94.5\% & 86.8\% & 6.3 \\
\bottomrule
\end{tabular}
\end{table}

The two HF H$_4$ rows are strongly depth-dominated: only $1$ to $3\%$ of
the cost is committed before the first elimination, and essentially all of it
goes into separating the leading pair. LiH, both H$_2$O geometries and the two
H$_4$ CISD rows are breadth-dominated, with $57$ to $94\%$ of the cost
already spent when every generator is still in play. What distinguishes them is
$\max_i\Delta_i/\mathrm{gap}$, the dynamic range the schedule has to cover:
$13.7$ and $18.9$ on the two depth-dominated H$_4$ HF rows against $1.2$ to
$3.4$ on H$_4$ CISD, LiH and H$_2$O equilibrium. Where the leading pair is
nearly tied relative to the spread of the field, depth dominates; where the
deficits are spread more evenly, breadth does. H$_2$O stretch, at $17.8$ and
$6.3$, does not fit that reading, so the correlation is suggestive rather than
established.

The two terms reward sharing very differently. In the breadth term a shared
context serves the entire pool, 92 or 140 generators at once, and M3's advantage
over M2 is correspondingly large. In the depth term it serves only the generators
that survive to the final radius. M3's breadth share is below M2's in every one of
the nine rows, which is what sharing is supposed to do: it makes the coarse sweep
cheaper relative to the rest.

\subsection{What decides the depth term}
\label{sec:fragmentation}

Two quantities decide it. Generator $i$'s own FC grouping is
$\mathcal M_{i,1},\ldots,\mathcal M_{i,K_i}$; write $F_{i\alpha}$ for the part of
$C_i$ landing in parent context $\mathcal M_\alpha$. The \emph{fragmentation
penalty}
\begin{equation}
  f_i=\left(\frac{\sum_{\alpha}\sigma_{i\alpha}\ \text{over parent contexts}}
                 {\sum_{k}\sigma_{i k}\ \text{over its own } \mathcal M_{i,k}}\right)^{2}
  \label{eq:fragmentation}
\end{equation}
is the factor by which the shared grouping makes one gradient more expensive to
resolve alone. It exceeds unity in every case here because
Eq.~\eqref{eq:single-gradient-shot} depends on $(\sum_\alpha\sigma_{i\alpha})^2$
rather than $\sum_\alpha\sigma_{i\alpha}^2$, so finer splitting raises the sum of
standard deviations --- but this is an \emph{observation}, not a theorem: the
parent grouping is not a refinement of the own grouping, since a parent context
can merge Paulis the own grouping separates, so a larger context count does not
by itself imply $f_i\ge1$. The second quantity is
$|\mathcal A|=|\{i:\Delta_i\le\mathrm{gap}\}|$, the number of generators still
active at the final radius; by construction $|\mathcal A|\ge2$.

In Table~\ref{tab:frag}, $K^{\text{own}}$ counts groups in the winner's own FC
grouping and depends only on $\supp(C_i)$, not the state; $K^{\text{parent}}$
counts parent contexts in which the winner has a \emph{non-zero} fragment
standard deviation, which is state-dependent --- hence $858$ against $887$, and
$875$ against $899$, on the H$_2$O HF and CISD rows for the same winner. Contexts
with $\sigma_{i\alpha}=0$ contribute nothing and are not counted.

\begin{table}[!htbp]
\centering
\scriptsize
\caption{Fragmentation of the winning gradient against the number of finalists. $K^{\text{own}}$ is the number of groups in the winner's own FC grouping; $K^{\text{parent}}$ is the number of parent contexts in which it has a non-zero fragment standard deviation, which is why it depends on the state. Both are defined in the text above.}
\label{tab:frag}
\begin{tabular}{lrrrrrr}
\toprule
Case & State & $K^{\text{own}}$ & $K^{\text{parent}}$ & $f$ & $|\mathcal A|$ & M2/M3 \\
\midrule
H$_4$, side $1.0$ & HF & 10 & 22 & 1.13 & 2 & 1.51 \\
H$_4$, side $1.0$ & CISD & 7 & 28 & 1.07 & 3 & 1.91 \\
H$_4$, side $2.0$ & HF & 10 & 22 & 1.19 & 3 & 1.92 \\
H$_4$, side $2.0$ & CISD & 7 & 28 & 1.22 & 2 & 2.60 \\
LiH, $R{=}3.0$ & HF & 58 & 520 & 3.49 & 3 & 1.79 \\
H$_2$O eq., $R{=}0.9572$ & HF & 57 & 858 & 2.54 & 2 & 1.13 \\
H$_2$O eq., $R{=}0.9572$ & CISD & 57 & 887 & 2.44 & 2 & 1.03 \\
H$_2$O str., $R{=}1.75$ & HF & 50 & 875 & 3.62 & 2 & 0.86 \\
H$_2$O str., $R{=}1.75$ & CISD & 50 & 899 & 2.94 & 2 & 1.01 \\
\bottomrule
\end{tabular}
\end{table}

The scaling looks asymmetric. $|\mathcal A|$ is set by how many generators are nearly tied at the top, which has nothing to do with pool size and is two or three in every case here. The fragmentation penalty is larger on the larger systems, because a larger universal support spreads each individual gradient over more parent contexts: 22 on H$_4$, 520 on LiH, roughly 875 on H$_2$O. If that persisted, M3's benefit in the depth term would have a ceiling while its cost in that term did not, and a crossover would follow.

\textbf{This is a hypothesis, not a demonstration.} $f$ is not monotone in system size --- LiH, the smaller system, gives $3.49$ against $2.54$ for H$_2$O at equilibrium --- four systems cannot separate size from geometry and state, and $f$ grows far more slowly than the raw context count ($3.6$ against $17.5$ on H$_2$O stretch), because most parent contexts hold a negligible sliver of any one gradient. The mechanism is real but weak, and nothing here fixes where or whether a crossover occurs. It is a statement about the planning bound, and within it about the depth term only; Sec.~\ref{sec:gamma} measures realised cost under a schedule, which the bound does not model, so the two do not compete.

\subsection{The finalists do not share Pauli terms}
\label{sec:overlap}

The premise of universal grouping is that gradient operators share Pauli products. Across the pool they do: each Pauli product in $\mathcal B_0$ appears in roughly four different gradient commutators.

\begin{table}[!htbp]
\centering
\scriptsize
\caption{Support sharing. Average reuse is $\sum_i|\supp(C_i)|/|\mathcal B_0|$. Finalist overlap is the Jaccard overlap $|\supp(C_{(1)})\cap\supp(C_{(2)})|/|\supp(C_{(1)})\cup\supp(C_{(2)})|$ between the supports of the two leading generators. Supports do not depend on the state, but the identity of the finalists does, so the two H$_4$ rows differ. The remaining four cases have not been evaluated.}
\label{tab:overlap}
\begin{tabular}{lrrrrr}
\toprule
Case & $|\mathcal B_0|$ & $\sum_i|\supp(C_i)|$ & Average reuse & Finalist overlap & M2/M3 \\
\midrule
H$_4$, side $1.0$, HF & 1,276 & 4,148 & 3.25 & 0.0\% & 1.51 \\
H$_4$, side $1.0$, CISD & 1,276 & 4,148 & 3.25 & 27.8\% & 1.91 \\
LiH, $R{=}3.0$, HF & 30,196 & 120,384 & 3.99 & 35.1\% & 1.79 \\
H$_2$O eq., $R{=}0.9572$, HF & 81,566 & 314,720 & 3.86 & 0.0\% & 1.13 \\
H$_2$O str., $R{=}1.75$, HF & 81,566 & 314,720 & 3.86 & 0.0\% & \textbf{0.86} \\
\bottomrule
\end{tabular}
\end{table}

The two columns tell opposite stories. Pool-wide reuse of about four is why M3's breadth term is efficient. But in three of the five cases the two finalists share \emph{no} Pauli products at all, so at the depth stage there is nothing to share and M3 pays the fragmentation penalty for nothing.

The overlap column does not order the M2/M3 column across systems. The two cases with substantial finalist overlap, LiH and H$_4$ CISD, are cases where M3 keeps a clear advantage, and the H$_2$O rows with none are where it disappears --- but H$_4$ HF has zero overlap and still gives $1.51$, well above both H$_2$O rows, and its fragmentation penalty is six times smaller than LiH's. Low finalist overlap is a plausible reason for M3 to lose its depth-stage advantage and is present where the advantage is smallest; five rows are not enough to make it \emph{the} explanation.

\subsection{Rebuilding the grouping once the field narrows}
\label{sec:hybrid}

If the fixed parent contexts are the problem at the depth stage, the obvious
remedy is to abandon them once the active set is small and build a fresh fully
commuting grouping over the survivors' union support. Phase one measures on the
parent contexts while more than a chosen number of generators remain active;
phase two regroups from scratch, using the same greedy first-fit routine and
ordering.

The earlier shots are not lost when this happens. They are taken in different
circuit settings, so no further shot can be added to the parent-context estimate
of a survivor, but that estimate remains unbiased at reconstruction variance
$V^{\mathrm{par}}_i=\sum_\alpha\sigma_{i\alpha}^2/n_\alpha$, and the regrouped
data give an independent one. The inverse-variance combination is unbiased at
$V^{-1}=(V^{\mathrm{par}})^{-1}+(V^{\mathrm{new}})^{-1}$, so phase two need only
supply the residual precision, arm by arm:
\begin{equation}
  \frac{1}{\epsilon_i^2}=\frac{1}{\epsilon^2}-\frac{1}{V^{\mathrm{par}}_i} ,
  \label{eq:reuse}
\end{equation}
imposed by scaling row $i$ of the new fragment standard deviations by
$1/\epsilon_i$ and solving Eq.~\eqref{eq:allocation} at unit epsilon. An arm whose
parent data already meet the target drops out of phase two entirely. Switch points
$|\mathcal A|\le2,3,4,6,10,20$ and every breakpoint of the trajectory were tested
on each case, both with this credit and with phase two charged the full cost of
reaching the final radii.

\begin{table}[!htbp]
\centering
\scriptsize
\caption{Staged regrouping against the better of M2 and M3, with the earlier shots credited by Eq.~\eqref{eq:reuse} and with them discarded. The switch point is the active-set size at which the parent contexts are abandoned.}
\label{tab:hybrid}
\begin{tabular}{lrrrrr}
\toprule
Case & Better of M2, M3 & Discard & Ratio & Reuse & Ratio \\
\midrule
H$_4$, side $1.0$, HF   & 566,951 & 502,858 & 0.89 & 494,591 & \textbf{0.87} \\
H$_4$, side $1.0$, CISD & 53,402  & 53,402 & 1.00 & 50,029 & \textbf{0.94} \\
H$_4$, side $2.0$, HF   & 220,388 & 182,709 & 0.83 & 181,899 & \textbf{0.83} \\
H$_4$, side $2.0$, CISD & 33,432  & 33,432 & 1.00 & 30,469 & \textbf{0.91} \\
LiH, $R{=}3.0$, HF      & 639,465 & 639,464 & 1.00 & 597,169 & \textbf{0.93} \\
\bottomrule
\end{tabular}
\end{table}

With the credit, staged regrouping pays on all five cases, by $6$ to $17\%$.
Without it, three of the five show no gain at all and the procedure degenerates to
M3. The saving is modest either way, and Table~\ref{tab:breadth-depth} says why:
regrouping can only improve the depth term, and on the systems where M3 struggles
the cost is dominated by breadth, which the remedy does not touch.

\textbf{The credit buys robustness more than it buys shots.} With the earlier
shots discarded, the penalty for switching at the wrong point is severe --- the
total runs from $1.02$ to $1.57$ of the better method across switch points on LiH
and on H$_4$ side $1.0$ CISD, and the best switch point is early on H$_4$ and late
on LiH. With the credit, the total is flat to within about $0.02$ over a wide
range and never exceeds $0.99$. The choice of switch point, which otherwise has to
be tuned per system with no rule that transfers, stops mattering.

\textbf{Where the leverage is.} It is at the first substantial elimination, not at
the final pair. On H$_4$ side $1.0$ HF the best switch is the first breakpoint,
where sixteen of twenty-six generators leave at once and regrouping the ten
survivors needs twelve contexts against fifty-three; the union support collapses
disproportionately because those sixteen zero-gradient generators hold $708$ of
the $1{,}276$ Pauli products in $\mathcal B_0$ ($55.5\%$) and their supports are
\emph{disjoint} from the ten non-zero ones. The same pattern holds on LiH, where
the best switch is $|\mathcal A|=34$ out of $92$.

This rules out an appealing alternative trigger. Since Sec.~\ref{sec:overlap}
shows the finalists often share no Pauli products at all, one might regroup at the
moment the active supports decouple. That criterion fires on only one of the five
cases, at $|\mathcal A|=2$, and there it gives $0.95$ against the $0.87$ available
from switching early; on the other four the active support never fully decouples
and it never fires. Overlap remains high --- $0.73$ to $0.78$ of the active support
is still shared --- exactly where the gain is largest, so high overlap and high
leverage coincide rather than trade off. The switch should be triggered by the
first large drop in $|\mathcal A|$, which is also where Part III symmetry
filtering would act, without measuring anything first.

\section{Attribution: the no-sharing baseline}
\label{sec:baseline}

The three methods are two binary choices, sharing and elimination, and they span
only three of the four combinations. The fourth is a baseline that does neither:
every generator is measured in its own fully commuting grouping, and every one of
them is resolved to the same common radius used by M1. It is reported here as a
diagnostic outside the method hierarchy, and it serves two purposes. It makes the
two mechanisms separately attributable, and it is a sanity bound, since measuring
every generator separately is always available as a fallback.

\begin{table}[!htbp]
\centering
\scriptsize
\caption{Improvement each mechanism buys against the no-sharing baseline of Table~\ref{tab:completed}. Sharing only is baseline/M1, elimination only baseline/M2, both baseline/M3. The product column is what ``both'' would be if the mechanisms acted independently; the interaction is the ratio of the two.}
\label{tab:attribution}
\begin{tabular}{llrrrrr}
\toprule
Case & State & Sharing only & Elimination only & Both & Product & Interaction \\
\midrule
LiH, $R{=}3.0$ & HF & 1.91 & 5.98 & 10.68 & 11.44 & 0.93 \\
H$_4$, side $1.0$ & HF & 2.16 & 5.67 & 8.55 & 12.25 & 0.70 \\
H$_4$, side $1.0$ & CISD & 2.19 & 1.27 & 2.43 & 2.77 & 0.88 \\
H$_4$, side $2.0$ & HF & 3.35 & 4.74 & 9.09 & 15.87 & 0.57 \\
H$_4$, side $2.0$ & CISD & 3.11 & 2.08 & 5.40 & 6.47 & 0.84 \\
H$_2$O eq., $R{=}0.9572$ & HF & 1.19 & 8.13 & 9.22 & 9.63 & 0.96 \\
H$_2$O eq., $R{=}0.9572$ & CISD & 1.18 & 10.44 & 10.78 & 12.30 & 0.88 \\
H$_2$O str., $R{=}1.75$ & HF & 1.14 & 252.97 & 218.59 & 288.50 & 0.76 \\
H$_2$O str., $R{=}1.75$ & CISD & 1.11 & 37.53 & 38.09 & 41.49 & 0.92 \\
\bottomrule
\end{tabular}
\end{table}

Two things follow, and both bear on how the project should be described.

\textbf{On the planning bounds the mechanisms are substitutes.} The interaction is
below unity in every one of the nine cases, between $0.57$ and $0.96$: applying
both buys less than the product of applying each alone, because elimination
removes precisely the situation that sharing exploits, many generators active at
once. \textbf{That reading holds only for the bound.} The bound resolves every M3
arm to the common radius, and Sec.~\ref{sec:spillover} shows that on LiH and
H$_2$O the shared allocation resolves the leader far more tightly than that; with
M3 eliminating on its actual precision the interaction exceeds unity on all five
of those rows. There the mechanisms compound. On H$_4$ they remain substitutes.

\textbf{Sharing contributes less than elimination, and less on the larger
systems.} On its own, universal grouping is worth a factor of $2.2$ to $3.4$ on
H$_4$, $1.9$ on LiH, and only $1.1$ to $1.2$ on H$_2$O. Elimination on its own is
worth between $1.27$ and $252.97$: the low end is the H$_4$ side $1.0$ CISD row,
where the gap is wide and there is little for elimination to save, and the high
end is the H$_2$O stretch HF row. The two ranges overlap, so ``sharing contributes
little'' should be read as \emph{relative to elimination} and not as an absolute
statement --- a factor of $3.4$ is a real saving on its own terms. What the table
does establish is the ordering of the two mechanisms and its direction: the
sharing factor is smallest on the largest system in this set, on all four H$_2$O
rows. Whether that is a size effect or a property of these particular molecules
cannot be settled from four systems.

\section{The M3 bound undercounts shared precision}
\label{sec:spillover}

Every M3 number so far resolves each active arm to the common radius $r$ and
removes arm $i$ once $2r<\Delta_i$. That is not what the allocation delivers. A
shared allocation is sized by its binding arm, so every other active arm is
resolved more tightly than required, and the parent contexts keep every shot
already taken, so the shots spent separating arms that have since left keep
sharpening the survivors. A run that eliminates on each arm's actual
reconstruction variance $V_i=\sum_\alpha\sigma_{i\alpha}^2/n_\alpha$ --- as any
real run computing its own variances would --- drops arm $i$ as soon as some $j$
satisfies
\begin{equation}
  |g_j|-z_\delta\sqrt{V_j}\ >\ |g_i|+z_\delta\sqrt{V_i} .
  \label{eq:actual}
\end{equation}
With $n_\alpha(r)=\max\!\left(n_\alpha^{\text{so far}},\,n_\alpha^{(A)}(r)\right)$
every $V_i$ is non-increasing as $r$ falls, so the dominance margin is monotone in
$r$ and each elimination radius is found by bisection. The trajectory is followed
event by event with no schedule, and cost is charged per context at its largest
demand, as in Eq.~\eqref{eq:m3-aggregate}. It is the $\gamma\to1$ limit of the
noiseless trial of Sec.~\ref{sec:finite-shot}, and agrees with it to within one
schedule step wherever that trial resolves (a test pins the agreement). M2 needs no
such column: each M2 arm is resolved to exactly the common radius, so its bound
already uses its actual precision. M1 and the baseline eliminate nothing.

\begin{table}[!htbp]
\centering
\scriptsize
\caption{M3 on the common-radius planning bound and eliminating on actual precision, Eq.~\eqref{eq:actual}. \emph{Breadth} is M3's breadth share from Table~\ref{tab:breadth-depth}. The interaction is that of Table~\ref{tab:attribution}, recomputed with the actual-precision M3; values above unity, where the two mechanisms compound rather than substitute, are set in bold.}
\label{tab:spillover}
\begin{tabular}{llrrrrrrrr}
\toprule
 & & & \multicolumn{3}{c}{M3} & \multicolumn{2}{c}{M2/M3} & \multicolumn{2}{c}{Interaction} \\
\cmidrule(lr){4-6}\cmidrule(lr){7-8}\cmidrule(lr){9-10}
Case & State & Breadth & Bound & Actual & Ratio & Bound & Actual & Bound & Actual \\
\midrule
LiH, $R{=}3.0$ & HF & 87.6\% & 639,465 & 333,272 & 0.52 & 1.78 & 3.42 & 0.93 & \textbf{1.79} \\
H$_4$, side $1.0$ & HF & 2.1\% & 566,951 & 569,020 & 1.00 & 1.51 & 1.50 & 0.70 & 0.70 \\
H$_4$, side $1.0$ & CISD & 76.7\% & 53,403 & 49,729 & 0.93 & 1.93 & 2.07 & 0.88 & 0.95 \\
H$_4$, side $2.0$ & HF & 0.8\% & 220,389 & 228,733 & 1.04 & 1.92 & 1.85 & 0.57 & 0.55 \\
H$_4$, side $2.0$ & CISD & 47.7\% & 33,432 & 31,553 & 0.94 & 2.60 & 2.75 & 0.84 & 0.89 \\
H$_2$O eq., $R{=}0.9572$ & HF & 85.5\% & 22,133,048 & 13,246,592 & 0.60 & 1.13 & 1.89 & 0.96 & \textbf{1.60} \\
H$_2$O eq., $R{=}0.9572$ & CISD & 81.2\% & 45,017,730,601 & 24,416,305,313 & 0.54 & 1.03 & 1.90 & 0.88 & \textbf{1.61} \\
H$_2$O str., $R{=}1.75$ & HF & 60.7\% & 20,163,542 & 14,581,169 & 0.72 & 0.86 & 1.19 & 0.76 & \textbf{1.05} \\
H$_2$O str., $R{=}1.75$ & CISD & 86.9\% & 165,495,967 & 122,188,158 & 0.74 & 1.02 & 1.38 & 0.92 & \textbf{1.25} \\
\bottomrule
\end{tabular}
\end{table}

\textbf{The bound is conservative on LiH and H$_2$O and accurate on H$_4$.}
Eliminating on actual precision costs M3 $0.52$ to $0.74$ of its planning bound on
all five LiH and H$_2$O rows, and $0.93$ to $1.04$ on all four H$_4$ rows. Where it
comes out slightly above the bound, the event-driven trajectory separates
eliminations that the bound lumps into one breakpoint, and the per-context maxima
of Eq.~\eqref{eq:m3-aggregate} then pay a little for the extra allocations, as in
Sec.~\ref{sec:m3-m1}.

\textbf{The saving comes from the leader.} A losing arm is removed once the leader
leads it by the sum of their two actual radii, where the bound charges both at $r$.
At the breadth breakpoint, where most of M3's cost on these systems sits, the
allocation is sized by the losers that are expensive to resolve; if the leader's
own fragments are cheap by comparison it is resolved far below the target. Its
standard error is then a fraction $\kappa$ of the target: $0.48$ on LiH and $0.56$ on
H$_2$O at equilibrium, against $0.94$ to $1.00$ on every H$_4$ row, where the
leader is itself binding. A binding loser is then released at $2/(1+\kappa)$ times
the bound's radius --- $1.35$ on LiH, about $1.8$ times fewer shots for the stage
that carries $88\%$ of the cost. On H$_2$O the losers are over-resolved as well
(median $\kappa$ of $0.70$ to $0.78$ against $1.00$ on LiH), which widens the
release radius further. This is why breadth share alone does not predict the
saving: H$_4$ side $1.0$ CISD is $77\%$ breadth but its leader binds, and it
gains $7\%$.

\textbf{Three conclusions change.} First, \emph{the M2-versus-M3 comparison
widens on every LiH and H$_2$O row}: to $3.42$ on LiH and $1.19$ to $1.90$ on
H$_2$O, including stretched H$_2$O HF, where the bound's $0.86$ becomes $1.19$ and
stays above unity under every grouping routine (Sec.~\ref{sec:routines}).
Second, \emph{the LiH finite-shot differential of Sec.~\ref{sec:gamma} is
accounted for}: the realised ratio of $3.04$ sits between the bound's $1.78$ and
the actual-precision $3.42$, noise costing M3 slightly more than it costs M2.
Third, \emph{the two mechanisms compound on LiH and H$_2$O}. The interaction
against the no-sharing baseline, below unity on every row at bound level, becomes
$1.05$ to $1.79$ on all five of those rows and stays at $0.55$ to $0.95$ on H$_4$.
Sharing spills precision onto arms the allocation was not sized for, and that
precision is exactly what lets elimination fire early. The substitution reading of
Sec.~\ref{sec:baseline} is a property of the common-radius bound, and holds for
M3 as it would actually run only where the leader binds.

\section{Finite-shot behaviour}
\label{sec:finite-shot}

Every number above is an oracle planning bound. The elimination rule reads exact
gradients, so it cannot misfire, and the family-wise confidence is asserted rather
than measured. This section replaces the exact gradients with estimates built from
a finite number of shots and runs the same elimination rule on those.

\textbf{Protocol.} A real run does not know the deficits in advance and cannot
jump to the exact breakpoints, so each trial follows a geometric radius schedule:
the common radius starts at $\max_i|g_i|$, wide enough that no arm is eliminated
for free, and is multiplied by $\gamma=0.5$ each round. At each round the
allocation of Eq.~\eqref{eq:allocation} is solved for the current active set and
radius, and each context's shot count is raised to that value if it is higher,
never lowered. Each context's fluctuation is carried as a running sum, each round
adding the draw for its new shots only, so $\widehat g_i$ is distributed as
$\mathcal N\!\left(g_i,\ \sum_\alpha\sigma_{i\alpha}^2/n_\alpha\right)$ at the
accumulated $n_\alpha$ and successive rounds are correlated as real accumulated
data would be. The elimination rule of Eq.~\eqref{eq:bai-elim} is then
applied. A trial ends when one arm remains or the radius falls below
$10^{-3}\max_i|g_i|$. Every case uses 200 trials at seed $0$; the trial count is
set so that the standard error on a mean cost is a few per cent, the resolution
the comparisons need. The sweeps below use 250, 300 and 500 trials: those tables
report \emph{rates} rather than costs, and the counts were raised to tighten the
Wilson intervals, which at 300 and 500 trials reach $\pm1.3\%$ and $\pm0.8\%$ at
zero observed errors. The starting radius and the floor are both taken
from the exact $\max_i|g_i|$ and are therefore oracle quantities, though mild ones:
a real run would take the starting radius from a cheap pilot estimate of the
largest gradient, and the floor only prevents a non-terminating run. This is a
weaker use of oracle knowledge than M1's exact gap (Sec.~\ref{sec:completed}),
but it is not none: the best-arm methods discover their \emph{stopping} point
from the data, not their starting point.

\textbf{The inflation is mostly the protocol, not noise.} Running the identical
trial loop with the noise draws set to zero isolates what the schedule and the
stopping rule cost on their own, at matched stopping rules, and the remainder is
what noise adds (Table~\ref{tab:shrink}). Schedule-only inflation is $1.74$ to
$3.48$ for M2 and $0.88$ to $2.85$ for M3, while noise multiplies it by only
$0.84$ to $1.40$ and $0.87$ to $2.67$. Several noise factors fall below unity, a
lucky fluctuation ending a run a round early. M3's noiseless trial cost falls
\emph{below} its planning bound on three rows, as low as $0.88$, because stopping
when one arm remains is a weaker criterion than resolving the leader to
$\mathrm{gap}/2$. The ``inflation'' column therefore mixes a weaker stopping
criterion, which lowers cost, with discretisation and noise, which raise it.

Both confounds are removed in Sec.~\ref{sec:gamma}, by sweeping the shrink factor
towards the exact limit and by matching the stopping rule to the planning
criterion. The result is that the inflation is overwhelmingly discretisation: at
$\gamma=0.95$ it falls to $1.03$--$1.35$ for M2 and $0.73$--$1.29$ for M3. The
M2/M3 differential visible in the table below survives on LiH and both H$_2$O CISD
rows, and Sec.~\ref{sec:spillover} shows it belongs to the bound.

The two deliberate approximations, the normal limit and independence between
generators, are listed in Sec.~\ref{sec:assumptions};
Sec.~\ref{sec:correlation} removes the second.

\textbf{Columns.} \emph{Bound} is the oracle planning cost of
Table~\ref{tab:completed}; \emph{Infl.} is the mean total shots over the 200
trials divided by it; \emph{Correct} is the fraction of trials returning the true
maximum-gradient generator. M2/M3 columns are ratios of the corresponding
entries, standard errors combined in quadrature.

\begin{table}[!htbp]
\centering
\scriptsize
\caption{Realised cost against the planning bound, 200 trials per row at seed $0$, all uncertainties one standard error of the mean. Every row selected the true maximiser in all 200 trials, a Wilson $95\%$ interval of $[0.981,1]$; the \emph{Correct} column is therefore constant and omitted. The non-adaptive baseline and M1 realise their bound exactly, having no schedule to overshoot, and are omitted.}
\label{tab:finite-shot}
\begin{tabular}{llrr@{\,$\pm$\,}lrr@{\,$\pm$\,}lrr@{\,$\pm$\,}l}
\toprule
& & \multicolumn{3}{c}{M2: BAI-FC-IG} & \multicolumn{3}{c}{M3: BAI-FC-UG} & \multicolumn{3}{c}{M2/M3} \\
\cmidrule(lr){3-5}\cmidrule(lr){6-8}\cmidrule(lr){9-11}
Case & State & Bound & \multicolumn{2}{c}{Infl.} & Bound & \multicolumn{2}{c}{Infl.} & Bound & \multicolumn{2}{c}{Realised} \\
\midrule
LiH, $R{=}3.0$ & HF & 1,140,583 & 3.16 & 0.06 & 638,423 & 1.40 & 0.03 & 1.79 & 4.03 & 0.12 \\
H$_4$, side $1.0$ & HF & 854,629 & 2.21 & 0.08 & 566,951 & 2.33 & 0.10 & 1.51 & 1.43 & 0.08 \\
H$_4$, side $1.0$ & CISD & 102,806 & 2.99 & 0.05 & 53,727 & 2.74 & 0.04 & 1.91 & 2.09 & 0.04 \\
H$_4$, side $2.0$ & HF & 422,746 & 2.33 & 0.05 & 220,389 & 2.48 & 0.05 & 1.92 & 1.81 & 0.05 \\
H$_4$, side $2.0$ & CISD & 86,857 & 2.77 & 0.06 & 33,432 & 2.95 & 0.09 & 2.60 & 2.43 & 0.09 \\
H$_2$O eq., $R{=}0.9572$ & HF & 25,096,606 & 3.02 & 0.06 & 22,130,852 & 1.53 & 0.04 & 1.13 & 2.24 & 0.07 \\
H$_2$O eq., $R{=}0.9572$ & CISD & 46,457,815,707 & 3.02 & 0.07 & 45,017,730,601 & 1.36 & 0.03 & 1.03 & 2.30 & 0.07 \\
H$_2$O str., $R{=}1.75$ & HF & 17,423,631 & 2.75 & 0.06 & 20,163,542 & 1.66 & 0.04 & \textbf{0.86} & 1.43 & 0.04 \\
H$_2$O str., $R{=}1.75$ & CISD & 168,305,383 & 2.95 & 0.07 & 165,495,967 & 2.07 & 0.05 & 1.02 & 1.45 & 0.05 \\
\bottomrule
\end{tabular}
\end{table}

Realised means are Bound $\times$ Infl.; the largest is $1.4\times10^{11}$ shots
for M2 on H$_2$O CISD at equilibrium and the smallest $9.9\times10^{4}$ for M3 on
H$_4$ side $2.0$ CISD.

The two H$_2$O CISD rows carry the largest bounds in the set, up to
$4.6\times10^{10}$ shots for M2 at equilibrium; the surrogate samples estimators
rather than outcomes, so shot totals of that size cost nothing to simulate.

The scatter is not negligible and is the reason the errors are quoted. Repeated runs of the same configuration at different seeds move a realised total by $2$ to $4\%$, which is comparable to some of the differences between adjacent rows and larger than the difference between the bound and realised M2/M3 ratios on the H$_4$ HF rows.

\subsection{The inflation is the radius schedule}
\label{sec:gamma}

The trials above use a geometric schedule of shrink factor $\gamma=0.5$. A real
run does not know the deficits $\Delta_i$ in advance, so it discovers them by
shrinking a common radius, and a coarse schedule overshoots every elimination
point. Sweeping $\gamma$ separates that overshoot from noise: as $\gamma\to1$ the
schedule approaches the exact-limit trajectory the planning bound assumes, so
whatever inflation survives is not discretisation. A second control replaces the
trial's stopping rule --- halt when one arm remains --- with the criterion the
bound assumes, halt at radius $\mathrm{gap}/2$; the default rule is the weaker of
the two and can stop before the bound's terminal radius.

\begin{table}[!htbp]
\centering
\scriptsize
\caption{Realised cost divided by planning bound, at the coarse schedule and near the exact limit, with the realised M2/M3 ratio against its bound-level value. 200 trials per entry (100 on LiH) at seed $0$, one-arm stopping rule; the matched rule changes no entry by more than $0.1$ except where noted in the text.}
\label{tab:gamma}
\begin{tabular}{lrrrrrrr}
\toprule
 & \multicolumn{2}{c}{M2 inflation} & \multicolumn{2}{c}{M3 inflation} & \multicolumn{2}{c}{realised M2/M3} & bound \\
\cmidrule(lr){2-3}\cmidrule(lr){4-5}\cmidrule(lr){6-7}
Case & $\gamma{=}0.5$ & $0.95$ & $\gamma{=}0.5$ & $0.95$ & $\gamma{=}0.5$ & $0.95$ & M2/M3 \\
\midrule
H$_4$, side $1.0$, HF   & 2.21 & 1.03 & 2.33 & 1.12 & 1.43 & 1.38 & 1.51 \\
H$_4$, side $1.0$, CISD & 2.99 & 1.35 & 2.76 & 1.29 & 2.09 & 2.02 & 1.93 \\
H$_4$, side $2.0$, HF   & 2.33 & 1.10 & 2.48 & 1.17 & 1.81 & 1.80 & 1.92 \\
H$_4$, side $2.0$, CISD & 2.77 & 1.21 & 2.95 & 1.19 & 2.43 & 2.64 & 2.60 \\
LiH, $R{=}3.0$, HF      & 3.01 & 1.27 & 1.40 & \textbf{0.74} & 3.84 & \textbf{3.04} & 1.78 \\
H$_2$O eq., CISD        & 3.02 & 1.29 & 1.36 & \textbf{0.73} & 2.30 & 1.81 & 1.03 \\
H$_2$O str., CISD       & 2.95 & 1.23 & 2.07 & \textbf{0.91} & 1.45 & 1.38 & 1.02 \\
\bottomrule
\end{tabular}
\end{table}

\textbf{Almost all of the inflation is the schedule.} At $\gamma=0.95$ M2 costs
$1.03$ to $1.35$ of its bound and M3 $0.73$ to $1.29$, against $2.2$ to $3.0$ at
$\gamma=0.5$. The coarse schedule was measuring its own overshoot. A deployment
would tune $\gamma$, and the cost of not tuning it is a factor of two to three.

\textbf{On H$_4$ the realised comparison reproduces the bound.} The realised
M2/M3 ratio at $\gamma=0.95$ is $1.38$, $2.02$, $1.80$ and $2.64$ against
bound-level $1.51$, $1.93$, $1.92$ and $2.60$, and the matched stopping rule moves
these to $1.43$, $2.03$, $1.87$ and $2.54$. The bound-level comparison of
Sec.~\ref{sec:completed} therefore carries over to finite shots on all four rows.

\textbf{LiH does not, and the reason is the bound, not the noise.} Its realised
ratio is $3.04$ under the one-arm rule and $3.12$ under the matched rule, against a
bound of $1.78$, and neither moves with $\gamma$: M2 inflates to $1.27$ while M3
falls to $0.74$ of its own bound. The same trial run with every noise draw set to
zero already costs M3 only $0.55$ of its bound at $\gamma=0.95$, with every arm
leaving at a radius above its exact $\Delta_i/2$; noise then raises M2 and M3 by
similar factors, $1.26$ and $1.38$. So the
trial is not beating the bound through lucky draws. It eliminates on each arm's
actual precision, which on a shared allocation is far better than the common radius
the bound assumes. Sec.~\ref{sec:spillover} quantifies this for every case; on LiH
it accounts for the whole differential.

\subsection{The comparison with M1 under finite shots}
\label{sec:m1-finite}

M1 and the no-sharing baseline are non-adaptive: they solve one allocation at one
radius and realise their planning bound exactly, with no schedule to overshoot.
The best-arm methods do not, so comparing realised cost against M1's bound is
sensitive to $\gamma$ in a way Table~\ref{tab:completed} is not.

\begin{table}[!htbp]
\centering
\scriptsize
\caption{M1 against the \emph{realised} cost of the best-arm methods at $\gamma=0.5$. M1's bound is its realised cost. Entries below unity, where the best-arm method is more expensive than static all-gradient estimation, are set in bold.}
\label{tab:m1-finite}
\begin{tabular}{llrrr}
\toprule
Case & State & M1 & M1 / M2 realised & M1 / M3 realised \\
\midrule
LiH, $R{=}3.0$ & HF & 3,563,904 & \textbf{0.99} & 3.99 \\
H$_4$, side $1.0$ & HF & 2,244,862 & 1.19 & 1.70 \\
H$_4$, side $1.0$ & CISD & 59,642 & \textbf{0.19} & \textbf{0.40} \\
H$_4$, side $2.0$ & HF & 598,102 & \textbf{0.61} & 1.09 \\
H$_4$, side $2.0$ & CISD & 58,068 & \textbf{0.24} & \textbf{0.59} \\
H$_2$O eq., $R{=}0.9572$ & HF & 172,147,028 & 2.27 & 5.07 \\
H$_2$O str., $R{=}1.75$ & HF & 3,864,803,902 & 80.7 & 115.3 \\
H$_2$O eq., $R{=}0.9572$ & CISD & 411,524,732,155 & 2.93 & 6.74 \\
H$_2$O str., $R{=}1.75$ & CISD & 5,715,348,740 & 11.5 & 16.7 \\
\bottomrule
\end{tabular}
\end{table}

At this schedule M3 loses to M1 on both H$_4$ CISD rows, which it won on the
bounds. \textbf{That reordering does not survive a finer schedule.} At
$\gamma=0.95$ M3 realises $39{,}840$ on H$_4$ side $2.0$ CISD against M1's
$58{,}068$, and with the matched stopping rule it realises $54{,}872$ on H$_4$
side $1.0$ CISD against M1's $59{,}642$: M3 beats M1 on all nine rows. The
inversions in the table are the $\gamma=0.5$ overshoot, not a property of finite
sampling, and the bound-level claim that M3 is cheaper than M1 stands under finite
shots as well.

What the table does show is where the margin is thin. Elimination pays where M1 is
expensive --- H$_2$O, where the factor stays between $2$ and $115$ --- and barely
pays where M1 is already cheap, which is the regime conclusion 4 identifies. M2,
which had already lost both CISD rows at bound level, also loses H$_4$ side $2.0$
HF and draws level on LiH.

\subsection{The one inversion does not survive}

On H$_2$O stretch with the Hartree--Fock reference the planning bounds give
M2/M3 $=0.86$, the only row in which M3 is more expensive than M2. Under finite
shots the same row gives $1.43\pm0.04$, and M3 wins. The inversion belongs to the
bound: the trial eliminates on each arm's actual precision, the bound does not,
and with M3 eliminating on actual precision the ratio is $1.08$ to $1.52$ across
all five grouping routines (Sec.~\ref{sec:routines}), where the bound itself moves
over $0.86$ to $1.05$. On actual precision M3 is the cheaper method on this row
under every routine tested.

\subsection{The confidence convention is far stricter than $\delta$, by an amount the data can only bound}
\label{sec:calibration}

No trial in the main finite-shot study of Table~\ref{tab:finite-shot} selected the
wrong generator, and sweeping the family-wise failure probability over nearly its
whole range barely changes that. Errors do appear elsewhere --- at $\delta=0.90$
and $0.99$ below, and at small $z_\delta$ in the lower block of Table~\ref{tab:calibration} --- so
the statement is about this sweep, not about every run in the report.

\textbf{How much stricter cannot be measured here.} With 300 trials and zero
errors a Wilson $95\%$ interval puts the rate anywhere between $0$ and $1.3\%$.
The $1.2\times10^{-5}$ quoted below is the analytic prediction of
Eq.~\eqref{eq:calibration}, not a measurement; resolving a rate that small needs
of order $10^{6}$ trials. What is established is that the realised rate is at most
$1.3\%$ where $0.05$ was asked for, about four times smaller, and that it barely
moves as $\delta$ is swept --- enough for the practical conclusion that $\delta$
is not the knob, not enough to say by how much the convention over-delivers.

\textbf{The same holds on two further cases.} Repeating the sweep over
$\delta\in\{0.05,0.5,0.9,0.99\}$ with 300 trials per entry, LiH returns no wrong
generator for either method at any $\delta$, and H$_4$ side $2.0$ HF none at
$\delta=0.05$ and at most $1.7\%$ [$0.7,3.8$] at $\delta=0.9$ --- a nominal
family-wise failure probability of ninety per cent delivering under two per cent.
The rate is set by the stopping radius and the gap, as below, and the confidence
level barely enters.

The reason is structural. Whether a run halts on one surviving arm or on the
radius, the leader can only be separated from the runner-up once the common
radius has come down to about $\mathrm{gap}/2$, at which point each estimate has
standard error $\mathrm{gap}/(2z_\delta)$. The leader's margin is one gap, so in
units of the standard deviation of the difference it is $z_\delta\sqrt2$, and
\begin{equation}
  P(\text{wrong selection})\approx 2\,\overline\Phi\!\left(z_\delta\sqrt2\right),
  \label{eq:calibration}
\end{equation}
where $\overline\Phi$ is the upper-tail standard normal function. The
single-competitor probability is $\overline\Phi(z_\delta\sqrt2)$; the factor of
two is a conservative allowance, not a derivation, covering the comparison being
between absolute values and more than one arm sitting near the top. Dropping it
halves every predicted rate and changes no conclusion. The Bonferroni term enters
only through $z_\delta$, and the stopping radius supplies most of the protection.

\begin{table}[!htbp]
\centering
\scriptsize
\caption{Confidence sweeps on H$_4$ side $1.0$ HF. Upper block: the stated $\delta$ against the realised selection-error rate, 300 trials per entry. Lower block: the confidence factor set directly, 500 trials per entry. Brackets are Wilson $95\%$ intervals. At these trial counts none of the upper rows distinguishes a rate of $10^{-5}$ from one of $10^{-2}$.}
\label{tab:calibration}
\begin{tabular}{llllrr}
\toprule
$\delta$ & $z_\delta$ & M2 error [95\% CI] & M3 error [95\% CI] & M2 shots & M3 shots \\
\midrule
0.05 & 3.10 & 0.0\% [0, 1.3] & 0.0\% [0, 1.3] & 1,891,179 & 1,324,843 \\
0.40 & 2.42 & 0.0\% [0, 1.3] & 0.0\% [0, 1.3] & 1,232,695 & 849,844 \\
0.90 & 2.11 & 0.7\% [0.2, 2.4] & 0.3\% [0.1, 1.9] & 1,128,410 & 691,975 \\
0.99 & 2.07 & 0.0\% [0, 1.3] & 0.3\% [0.1, 1.9] & 969,099 & 741,804 \\
\midrule
--- & 1.39 & 12.2\% [9.6, 15.4] & 10.6\% [8.2, 13.6] & 1,586,040 & 989,854 \\
--- & 1.80 & 1.8\% [1.0, 3.4] & 2.0\% [1.1, 3.6] & 1,113,341 & 819,684 \\
--- & 2.10 & 0.4\% [0.1, 1.5] & 0.0\% [0, 0.8] & 1,211,688 & 685,567 \\
--- & 3.10 & 0.0\% [0, 0.8] & 0.0\% [0, 0.8] & 2,004,737 & 1,335,779 \\
\bottomrule
\end{tabular}
\end{table}

Eq.~\eqref{eq:calibration} predicts $0.049$, $0.011$, $0.003$ and
$1.2\times10^{-5}$ at $z_\delta=1.39$, $1.80$, $2.10$ and $3.10$. It accounts only for the final comparison and ignores
generators wrongly eliminated earlier in the run, so it underestimates the error
at small $z_\delta$: it predicts $5\%$ at $z_\delta=1.39$ where the
simulation gives ten to twelve. The cost is also not monotone. Below
$z_\delta\approx1.8$ it rises again, because bad eliminations prolong the run.
There is an optimum near $z_\delta\approx1.8$ to $2.1$, worth about $1.7$ in
realised cost at an error rate whose $95\%$ interval reaches $3.6\%$ at
worst --- a real saving, but not the factor of five the bound suggested. Both
conventions are selectable in the implementation; the tables in this report retain
the Bonferroni one. The sweep was run on one case, H$_4$ side $1.0$ with the
Hartree--Fock reference, and the optimum has not been checked on any other, so
$z_\delta\approx1.8$ to $2.1$ should be treated as a single-system observation
rather than a recommended default.

\subsection{Good-enough selection}
\label{sec:good-enough}

A run does not have to identify the maximiser; it is often enough to return a
generator that is not much worse. Fix a tolerance $\tau>0$ --- written $\tau$ to
keep it distinct from the deficits $\Delta_i$ --- and stop at the first round $t$
at which every surviving candidate is certified within $\tau$ of the best,
\begin{equation}
  \max_{j\in\mathcal A_t}\left(|\widehat g_j|-r_j\right)
  \ \ge\
  \max_{i\in\mathcal A_t}\left(|\widehat g_i|+r_i\right)-\tau ,
  \label{eq:good-enough}
\end{equation}
returning $\argmax_{i\in\mathcal A_t}|\widehat g_i|$. The run is counted a success
if the returned generator's true absolute gradient is within $\tau$ of the true
maximum.

On H$_4$ side $1.0$ HF at $\tau=\mathrm{gap}/4=5.38\times10^{-3}$, over the same
200 trials at seed $0$ as Table~\ref{tab:finite-shot}, M2 falls from
$1{,}889{,}539\pm64{,}528$ to $1{,}398{,}814\pm32{,}966$ shots and M3 from
$1{,}320{,}830\pm57{,}944$ to $904{,}382\pm20{,}201$, savings of $26\%$ and $31\%$.
In all 200 trials the generator returned was still the exact maximiser. Over 300
trials each, the same tolerance saves $18\%$ (M2) and $24\%$ (M3) on H$_4$ side
$2.0$ HF but only $2$ to $3\%$ on LiH, in every trial returning a generator
within $\tau$ of the best. The saving tracks how crowded the top of the spectrum
is: $\tau=\mathrm{gap}/4$ spares the final tightening, which is most of the cost
where the leading gradients are close together and almost none of it where the
gap is wide. Note that $\tau=\mathrm{gap}/4$ uses the exact gap, so this inherits
the oracle-knowledge caveat of Sec.~\ref{sec:completed}; a deployment would set
$\tau$ from a chemical accuracy target instead.

\section{Parameter sensitivity}
\label{sec:sensitivity}

\subsection{Removing the two free parameters}

Two arbitrary constants have been removed. M1 is evaluated at a target radius of
$\mathrm{gap}/2$, the loosest common radius at which a non-adaptive estimate
separates the winner from the field, so that all three methods answer the same
question. The BAI methods are evaluated in the schedule-free limit of
Sec.~\ref{sec:proxy} rather than on a geometric schedule, because a coarse
schedule overshoots each elimination point and overpays.

Neither change is cosmetic. Under $\gamma=0.5$, M3 costs more than M1 on three
H$_4$ rows. In the exact limit it does not, on any row
(Sec.~\ref{sec:completed}), so that ordering is a property of the limit rather
than of the method; under a coarse schedule it can fail, and
Sec.~\ref{sec:m1-finite} shows the same failure appearing and then disappearing
under finite shots as $\gamma$ is refined.

\begin{table}[!htbp]
\centering
\scriptsize
\caption{Decomposition of the finite-shot cost at matched stopping rules. \emph{Schedule} is the identical trial loop with the noise draws set to zero, divided by the planning bound; \emph{noise} is the realised mean divided by that; their product is the realised inflation of Table~\ref{tab:finite-shot}. The last three columns track the M2/M3 ratio through the same three stages.}
\label{tab:shrink}
\begin{tabular}{llrrrrrrrr}
\toprule
& & \multicolumn{2}{c}{M2} & \multicolumn{2}{c}{M3} & \multicolumn{3}{c}{M2/M3} \\
\cmidrule(lr){3-4}\cmidrule(lr){5-6}\cmidrule(lr){7-9}
Case & State & Schedule & Noise & Schedule & Noise & Bound & Schedule & Realised \\
\midrule
LiH, $R{=}3.0$ & HF & 2.25 & 1.40 & \textbf{0.88} & 1.59 & 1.79 & 4.56 & 4.03 \\
H$_4$, side $1.0$ & HF & 1.74 & 1.27 & 1.62 & 1.44 & 1.51 & 1.62 & 1.43 \\
H$_4$, side $1.0$ & CISD & 3.48 & 0.86 & 2.16 & 1.27 & 1.91 & 3.09 & 2.09 \\
H$_4$, side $2.0$ & HF & 2.75 & 0.85 & 2.85 & 0.87 & 1.92 & 1.85 & 1.81 \\
H$_4$, side $2.0$ & CISD & 2.28 & 1.21 & 1.11 & 2.67 & 2.60 & 5.36 & 2.43 \\
H$_2$O eq., $R{=}0.9572$ & HF & 3.37 & 0.90 & \textbf{0.88} & 1.74 & 1.13 & 4.34 & 2.24 \\
H$_2$O str., $R{=}1.75$ & HF & 3.26 & 0.84 & 1.91 & 0.87 & \textbf{0.86} & 1.48 & 1.43 \\
\bottomrule
\end{tabular}
\end{table}

Bold marks rows where the noiseless trial costs \emph{less} than the planning
bound. Under the planning stopping rule instead, the same schedule overpays by
$1.64$ to $3.80$ and moves the M2/M3 ratio by only $-4\%$ to $+20\%$; the
contrast with the Schedule column measures how much the stopping rule, rather
than the shrink factor, is doing. Refining the schedule converges onto the exact
limit as it must:Refining the schedule converges onto the exact limit as it must: on H$_4$ side
$1.0$ HF, M2 falls $1{,}494{,}274\to1{,}267{,}755\to1{,}080{,}113\to885{,}115$
and M3 $928{,}366\to872{,}427\to712{,}219\to584{,}142$ at
$\gamma=0.5,0.7,0.85,0.97$, against exact limits of $854{,}629$ and $566{,}951$.
\texttt{verify\_part1.py} asserts this convergence.

For M2 the limit is the closed form
\begin{equation}
  M_{\text{M2}}
  =\sum_i\left(\frac{z_\delta\sum_\alpha\sigma_{i\alpha}}{\Delta_i/2}\right)^{2},
  \label{eq:m2-exact}
\end{equation}
with $\Delta_i$ replaced by the gap for the leader. For M3 the trajectory has one
breakpoint per distinct $\Delta_i$ and the allocation of Eq.~\eqref{eq:allocation}
is evaluated at each.

\subsection{The confidence level does not change the comparisons}

Every planning bound scales with $z_\delta^2$, so every ratio in
Table~\ref{tab:completed} is invariant under the family-wise failure probability.
Verified for $\delta\in\{0.2,0.05,0.01,0.001\}$: the H$_4$ side $1.0$, HF ratios stay
at $2.63$ and $1.51$ to better than one part in $10^3$, while the absolute costs
move by about a factor of two.

\textbf{This invariance is a property of the bounds and does not extend to
realised cost.} Once a schedule and finite shots are introduced, $z_\delta$ changes
where each elimination fires and therefore how much each method overshoots, and
the two methods are not affected equally: Table~\ref{tab:calibration} gives
realised M2/M3 ratios of $1.60$, $1.36$, $1.77$ and $1.50$ at
$z_\delta=1.39$, $1.80$, $2.10$ and $3.10$ on one H$_4$ row --- a spread of $1.3$,
with no monotone pattern. What $\delta$ does to the realised error rate is treated
in Sec.~\ref{sec:calibration}.

\section{Routine independence}
\label{sec:routines}

Every number above comes from one first-fit grouping with one insertion order.
First-fit is deterministic given an order, but the order is a free convention, so
the conclusions have to be tested against it. The insertion order is a parameter:
three deterministic conventions --- by decreasing weight (\texttt{weight}, the
default used everywhere else), by increasing weight (\texttt{reverse-weight}) and
lexicographic --- together with a \texttt{random} family under seeds $0,1,2,3$.
The H$_4$ rows use all seven, LiH the three deterministic orderings and
\texttt{random} seed $0$, the H$_2$O rows the three deterministic orderings and
\texttt{random} seeds $0,1$. The ordering is applied to every grouping the
methods use --- the parent contexts of M1 and M3 and the individual groupings of M2
and the baseline --- so each row recomputes all four constructions from scratch.
Every combination is recorded per case in \texttt{*\_grouping\_routines.csv}.

\begin{table}[!htbp]
\centering
\scriptsize
\caption{Spread of the method ratios across groupings of the same problem, each recomputing the baseline and all three methods. The last column records whether the ordering of the methods is the same under every grouping tested.}
\label{tab:routines}
\begin{tabular}{lrrrrc}
\toprule
Case & Contexts & M1/M2 & M2/M3 & M1/M3 & Ordering fixed \\
\midrule
H$_4$, side $1.0$, HF   & 49--61 & 1.91--2.63 & 1.51--1.97 & 3.35--4.55 & yes \\
H$_4$, side $1.0$, CISD & 49--61 & 0.40--0.58 & 1.93--2.78 & 1.10--1.15 & yes \\
H$_4$, side $2.0$, HF   & 49--61 & 1.21--1.54 & 1.92--2.56 & 2.61--3.33 & yes \\
H$_4$, side $2.0$, CISD & 49--61 & 0.44--0.72 & 2.56--3.92 & 1.71--1.91 & yes \\
LiH, $R{=}3.0$, HF      & 980--1,127 & 2.58--3.12 & 1.79--2.11 & 5.22--5.58 & yes \\
H$_2$O eq., HF          & 2,366--2,849 & 5.63--6.86 & 1.13--1.36 & 7.65--7.78 & yes \\
H$_2$O eq., CISD        & 2,366--2,849 & 7.20--8.86 & 1.03--1.34 & 9.14--9.63 & yes \\
H$_2$O str., HF         & 2,356--2,870 & 181--250 & \textbf{0.86--1.05} & 172--262 & \textbf{no} \\
H$_2$O str., CISD       & 2,356--2,870 & 27.4--38.0 & \textbf{0.94--1.24} & 33.9--35.7 & \textbf{no} \\
\bottomrule
\end{tabular}
\end{table}

On seven of the nine rows the ratios move by at most a factor of $1.64$ and the
ordering of the methods is identical under every grouping tested. Those rows are
routine-independent for the quantities in the table.

\textbf{The two stretched H$_2$O rows are the exceptions, and the only ones.} Their
M2/M3 ratios run over $0.86$ to $1.05$ (HF) and $0.94$ to $1.24$ (CISD), spanning
unity: which of the two best-arm methods is cheaper on the planning bound is itself
a function of an arbitrary convention. At equilibrium the same five orderings give
$1.13$ to $1.36$ and $1.03$ to $1.34$, entirely above unity, so both equilibrium
rows \emph{are} resolved in M3's favour; on each, the value quoted in
Table~\ref{tab:completed} is the least favourable of the five. The two geometries
differ here for the same reason their Pauli supports differ
(Sec.~\ref{sec:structure}). The M1 versus M3 comparison is never in doubt on any
H$_2$O row, at $7.7$ to $262$ across the sweep. The stretched CISD row's canonical-routine
ratio reads $1.02$ in the sweep against $1.01$ in Table~\ref{tab:completed} for the
reason given in Sec.~\ref{sec:proxy}.

\textbf{On actual precision both stretched rows are resolved.} Repeating the
sweep with M3 eliminating on each arm's actual precision
(Sec.~\ref{sec:spillover}) gives M2/M3 of $1.08$ to $1.52$ on stretched H$_2$O HF
and $1.35$ to $2.36$ on stretched H$_2$O CISD, across the same five orderings: M3
is the cheaper method under every one. The irresolution above is a property of
the common-radius bound, which charges M3 for precision its allocation delivers
anyway; it is not a property of the method.

\section{Within-context correlation}
\label{sec:correlation}

Assumption (iii) of Sec.~\ref{sec:assumptions} is removed here: two generators
read off the same shots in a shared context, so their estimates are correlated,
and modelling that correlation makes the simulation exact for M1 and M3 as well.

\textbf{The context covariance.} For each context $\alpha$, let
$c_{i\alpha}=\left(F_{i\alpha}-\expect{F_{i\alpha}}\right)\ket{\psi}$ be the
centred fragment vector. The covariance is the Gram matrix
\begin{equation}
  \Cov(F_{i\alpha},F_{j\alpha})=\Re\left\langle c_{i\alpha}\middle| c_{j\alpha}\right\rangle ,
  \label{eq:gram}
\end{equation}
which is positive semidefinite by construction and whose diagonal is exactly
$\sigma_{i\alpha}^2$. Both properties are load-bearing and neither is automatic.
The polarisation identity
$\Cov(A,B)=\tfrac12[\Var(A{+}B)-\Var(A)-\Var(B)]$ is equivalent in exact
arithmetic, but a parent context is dominated by symmetry-forbidden fragments, for
which $\ket{\psi}$ is an eigenstate of $F_{i\alpha}$ and the variance is zero; there
the identity subtracts three quantities that are individually at the floating-point
noise floor. Computing each variance as $\expect{F^2}-\expect{F}^2$ compounds the
problem, since that form also cancels catastrophically for a near-eigenstate
fragment. Together the two produced correlations as large as $136$ and context
matrices with relative eigenvalues as low as $-6\times10^{264}$. Eq.~\eqref{eq:gram}
with $\Var(F)=\|c\|^2$ removes both cancellations: measured over all $53$ parent
contexts of the H$_4$ rows the worst relative eigenvalue is $-7\times10^{-16}$, no
pair violates Cauchy--Schwarz, and every diagonal matches $\sigma_{i\alpha}^2$ to
$10^{-15}$. The test suite pins all three properties (Sec.~\ref{sec:tests}); the
consequence of losing them is recorded at the end of this section.

\subsection{The covariance-aware elimination rule}

Where the \emph{marginal} rule of Eq.~\eqref{eq:bai-elim} compares
$|\widehat g_i|$ and $|\widehat g_j|$ at their individual radii, the
\emph{pairwise} rule compares each pair at the radius of the combination the
comparison actually depends on. Because the comparison is between absolute
values, that combination is $\widehat g_i-s_{ij}\widehat g_j$ with
$s_{ij}=\operatorname{sign}(\widehat g_i\widehat g_j)$, so $i$ is eliminated when
some active $j$ satisfies
\begin{equation}
  |\widehat g_j|-|\widehat g_i| \ > \ z_\delta\sqrt{V_{ij}},
  \qquad
  V_{ij}=\sum_\alpha\frac{1}{n_\alpha}
         \Var\!\left(F_{i\alpha}-s_{ij}F_{j\alpha}\right),
  \label{eq:pairwise}
\end{equation}
with the same $z_\delta$ of Eq.~\eqref{eq:zdelta}. When the two fragments are
correlated in the direction that helps, $V_{ij}$ is below
$\Var(\widehat g_i)+\Var(\widehat g_j)$ and the rule fires earlier.

The sign factor is required. Using $\Var(\widehat g_i-\widehat g_j)$
unconditionally is not only the wrong combination for an opposite-sign pair; it is
not well defined, because the sign of a gradient depends on the arbitrary phase of
the Hartree--Fock orbital it excites from. Only $s_{ij}\Cov(\widehat g_i,\widehat g_j)$
is phase-invariant, and Eq.~\eqref{eq:pairwise} uses exactly that.

\begin{table}[!htbp]
\centering
\scriptsize
\caption{M3 under correct joint noise, 250 trials per row at seed $0$. Uncertainties are one standard error; brackets are Wilson $95\%$ intervals on the correct-selection rate.}
\label{tab:correlation}
\begin{tabular}{llr@{\,$\pm$\,}llr}
\toprule
Case & Noise model, rule & \multicolumn{2}{c}{Shots} & Correct [95\% CI] & vs indep. \\
\midrule
H$_4$, side $1.0$, HF & independent, marginal & 1,334,743 & 51,258 & 100.0\% [98.5, 100] & --- \\
 & correlated, marginal & 1,323,818 & 51,343 & 100.0\% [98.5, 100] & 0.99 \\
 & correlated, pairwise & 685,298 & 23,477 & 100.0\% [98.5, 100] & 0.51 \\
H$_4$, side $1.0$, CISD & independent, marginal & 148,219 & 1,989 & 100.0\% [98.5, 100] & --- \\
 & correlated, marginal & 151,312 & 1,604 & 100.0\% [98.5, 100] & 1.02 \\
 & correlated, pairwise & 93,530 & 2,619 & 100.0\% [98.5, 100] & 0.63 \\
H$_4$, side $2.0$, HF & independent, marginal & 544,228 & 10,397 & 100.0\% [98.5, 100] & --- \\
 & correlated, marginal & 522,022 & 11,231 & 100.0\% [98.5, 100] & 0.96 \\
 & correlated, pairwise & 366,027 & 14,242 & 100.0\% [98.5, 100] & 0.67 \\
\bottomrule
\end{tabular}
\end{table}

Modelling the correlation changes the cost by $1$ to $4\%$ on all three rows. The
covariance-aware rule then cuts M3's realised cost by $33$ to $49\%$ at $100.0\%$
[$98.5,100$] correct selection throughout. This is the mechanism
Part II~\cite{PartII} proposes for the decisive stage, appearing as a direct
reduction in shots rather than as the outcome of an optimisation over coefficient
splitting.

\textbf{Why the construction of Eq.~\eqref{eq:gram} matters.} With the covariance
built from the polarisation identity instead, H$_4$ side $2.0$ HF returned
$75.6\%$ [$69.9,80.5$] correct under the marginal rule and $51.2\%$ [$45.0,57.3$]
under the pairwise rule, at eleven to seventeen times the cost. The mechanism is
specific: an indefinite context matrix makes the accumulated $V_{ij}$ of
Eq.~\eqref{eq:pairwise} negative, the implementation clamps it at zero before
taking the square root, and the arm is then eliminated on any lead whatever. On
that case every trial hit a negative $V_{ij}$; in a representative one the true
leader was removed at round three against a lead of $1.8\times10^{-2}$ at radius
zero, from $\Var_i=6.4\times10^{-4}$, $\Var_j=4.7\times10^{-4}$ and a reported
covariance of $2.1\times10^{-3}$, an implied correlation of $3.8$. The failure is
loudest where the runner-up is exactly degenerate --- \texttt{D 0,3 $\to$ 4,7} and
\texttt{D 1,2 $\to$ 5,6} both at $|g|=0.298387$, so $|\mathcal A|=3$ --- and the
gap is only $5\%$ of $|g_{(1)}|$. It is recorded here because a negative variance
silently clamped to zero is not a visible failure, and because the two properties
that exclude it, positive semidefiniteness and an exact diagonal, are cheap to
test and were not being tested.

\textbf{A diagnostic, defined once.} For an active pair write
\begin{equation}
  R_{ij}=\frac{\Var\!\left(\widehat g_i-s_{ij}\widehat g_j\right)}
              {\Var(\widehat g_i)+\Var(\widehat g_j)}
        =1-\frac{2\,s_{ij}\Cov(\widehat g_i,\widehat g_j)}
                {\Var(\widehat g_i)+\Var(\widehat g_j)} ,
  \label{eq:Rij}
\end{equation}
the factor by which the decisive comparison is harder ($R_{ij}>1$) or easier
($R_{ij}<1$) than the marginal radii assume. Because $s_{ij}$ and
$\Cov$ change sign together under an orbital phase flip, $R_{ij}$ is
phase-invariant. Over the three cases it is $0.929$, $1.084$ and $1.168$.

\section{Conclusions from the current data}
\label{sec:conclusions}

Table~\ref{tab:evidence} gives the evidence level; items resting on one system or
on bounds alone should not be generalised without the work in
Sec.~\ref{sec:remaining}.

\begin{enumerate}[leftmargin=2em]
  \item The four constructions are implemented as separate modules saving their FC groups, fragment variances and elimination histories, evaluated at one criterion with no free schedule or precision parameter. The validation gate passes on all nine cases to $4\times10^{-15}$ Ha.
  \item \textbf{M3 is cheaper than M1 in every case}, by $1.1$ to $192$ on the planning bounds and on realised cost in all nine simulated rows. \textbf{This is not a theorem}: under the sum-of-maxima accounting of Eq.~\eqref{eq:m3-aggregate} small problems with a nearly tied runner-up give $M_{\mathrm{M3}}/M_{\mathrm{M1}}$ up to $1.09$ (Sec.~\ref{sec:m3-m1}). Every M1 column assumes M1 is handed the exact gap.
  \item \textbf{M3 is cheaper than M2 on every row once it eliminates on actual precision.} On the planning bound M3 wins eight rows of nine, and both equilibrium H$_2$O rows are robust to the grouping routine ($1.13$--$1.36$ and $1.03$--$1.34$); both stretched rows are not, their bound-level ratios spanning unity ($0.86$--$1.05$ and $0.94$--$1.24$). With M3 eliminating on actual precision those two become $1.08$--$1.52$ and $1.35$--$2.36$ across the same routines, and the finite-shot trial agrees ($1.43$ and $1.45$ at $\gamma=0.5$). The bound-level irresolution is the bound's (Secs.~\ref{sec:routines} and \ref{sec:spillover}).
  \item \textbf{M2 is not always cheaper than M1}: seven bound-level rows of nine, five of nine on realised cost at $\gamma=0.5$. The counterexamples are wide-gap rows where static estimation is already cheap, but LiH has a comparably wide gap and does not invert, so the gap alone is not the explanation.
  \item \textbf{On the planning bounds the two mechanisms are substitutes; as M3 would actually run, they compound on LiH and H$_2$O.} At bound level the interaction against the no-sharing baseline is $0.57$ to $0.96$ in all nine cases. With M3 eliminating on its actual precision it is $1.05$ to $1.79$ on all five LiH and H$_2$O rows and $0.55$ to $0.95$ on H$_4$ (Sec.~\ref{sec:spillover}).
  \item \textbf{Elimination dominates in seven rows of nine.} Elimination alone is worth $1.27$ to $252.97$, sharing alone $1.11$ to $3.35$; sharing wins on both H$_4$ CISD rows. The sharing factor is smallest on all four H$_2$O rows, which with four systems is a correlation with system, not a trend with size, and refers to bounds only.
  \item \textbf{The finite-shot inflation is the radius schedule, not noise.} At the coarse schedule $\gamma=0.5$ both methods cost $2.2$ to $3.0$ times their bounds. Taking $\gamma$ to $0.95$ collapses that to $1.03$--$1.35$ for M2 and $0.73$--$1.29$ for M3, so what the coarse schedule measures is its own overshoot of each elimination point. The realised M2/M3 ratio then tracks the bound-level ratio on all four H$_4$ rows under both stopping rules. On LiH and both H$_2$O CISD rows it does not: $3.04$, $1.81$ and $1.38$ against bounds of $1.78$, $1.03$ and $1.02$. The trial eliminates on actual precision and the bound does not, and the actual-precision M3 of Sec.~\ref{sec:spillover} puts those ratios at $3.42$, $1.90$ and $1.38$, so the differential is the bound's, not the noise's.
  \item \textbf{The confidence convention does not control the selection-error rate.} Sweeping $\delta$ from $0.05$ to $0.99$ moves the realised rate only within $0$--$1.7\%$ on three cases, and at $\delta=0.05$ 300 trials bound it below $1.3\%$ on each. Inverting Eq.~\eqref{eq:calibration} for $5\%$ gives $10$ to $12\%$ in simulation, and an optimum sits near $z_\delta\approx1.8$--$2.1$, both on H$_4$ side $1.0$ HF. Good-enough stopping at $\tau=\mathrm{gap}/4$ saves $18$ to $31\%$ where the top of the spectrum is crowded and $2$ to $3\%$ on LiH, always returning a generator within $\tau$.
  \item \textbf{Conclusions 2 to 6 are routine-independent at bound level on seven of the nine rows}: all four H$_4$ rows, both H$_2$O equilibrium rows and LiH, each recomputed under five to seven groupings of every kind the methods use. Ratios move by at most $1.64$ and no ordering changes. The two stretched H$_2$O rows are the exceptions (conclusion 3). The finite-shot and correlated studies are unswept.
  \item \textbf{Within-context correlation is a saving on all three cases tested.} Modelling it changes cost by $1$ to $4\%$, and the covariance-aware pairwise rule then cuts M3's realised cost by $33$ to $49\%$ at $100.0\%$ [$98.5,100$] correct selection throughout. The saving requires a context covariance that is positive semidefinite with an exact diagonal (Sec.~\ref{sec:correlation}).
  \item \textbf{Regrouping once the field narrows is worth $6$ to $17\%$}, on all five cases tested, provided the shots already taken are credited by inverse-variance combination rather than discarded. Discarding them leaves three of the five cases with no gain at all. The credit matters less for the size of the saving than for its robustness: with it the total is flat to within $0.02$ across a wide range of switch points, without it the penalty for switching late reaches $1.57$ (Sec.~\ref{sec:hybrid}).
  \item \textbf{The M3 planning bound is conservative on the larger systems.} Eliminating on actual precision, M3 costs $0.52$ to $0.74$ of its bound on all five LiH and H$_2$O rows and $0.93$ to $1.04$ on H$_4$. The saving comes from the leader: where its fragments are cheap relative to the losers that size the breadth allocation, it is resolved to $0.48$--$0.79$ of the target, and losers leave at up to $1.35$ times the bound's radius (Sec.~\ref{sec:spillover}). The finite-shot trial eliminates the same way, which is why M3 realises below its bound on LiH and H$_2$O.
  \item All numbers are single-step, single-state diagnostics counting shots only --- not a full ADAPT trajectory, and not circuit-level cost, where M2 and the baseline need $3.5$ to $6.1$ times as many distinct settings (Table~\ref{tab:settings}).
\end{enumerate}

\begin{table}[!htbp]
\centering
\scriptsize
\caption{Evidence behind each conclusion. ``Bound'' is the oracle planning model, ``finite-shot'' the surrogate of Sec.~\ref{sec:finite-shot}, ``swept'' the number of rows on which the claim survives every grouping of Sec.~\ref{sec:routines}.}
\label{tab:evidence}
\begin{tabular}{clllc}
\toprule
\# & Claim & Bound & Finite-shot & Swept \\
\midrule
2  & M3 $<$ M1                        & 9 of 9  & 9 of 9 at $\gamma=0.95$ & 9 of 9 \\
3  & M3 $<$ M2                        & 8 of 9; 9 of 9 actual & 9 of 9 & 7 of 9; 9 of 9 actual \\
4  & M2 $<$ M1                        & 7 of 9  & 5 of 9 at $\gamma=0.5$ & 9 of 9 \\
5  & Substitutes                      & 9 of 9; 4 of 9 actual & ---    & no \\
6  & Elimination dominates            & 7 of 9  & ---    & 9 of 9 \\
7  & Inflation is the schedule        & ---     & 7 cases, 2--5 values of $\gamma$ & no \\
8  & $\delta$ does not control error  & ---     & 3 cases, 300--500 trials & no \\
10 & Covariance-aware rule            & ---     & 3 cases, 250 trials & no \\
11 & Regrouping with reuse            & 5 cases & ---    & no \\
12 & M3 bound conservative            & 9 cases & 2 cases, noiseless trial & no \\
13 & Distinct-setting counts          & 4 systems & ---  & no \\
\bottomrule
\end{tabular}
\end{table}

\section{Interface with Parts II and III}

Part II improves M3 by using measured variances and covariances to optimise overlapping universal operators and coefficient splitting. Sec.~\ref{sec:breadth-depth} gives that work a sharper target, but Sec.~\ref{sec:overlap} also raises a caution that should be settled before it is built.

The mechanism Part II~\cite{PartII} proposes for the decisive stage is a pairwise-variance identity: best-arm identification only requires the pairwise differences to be resolved, and, with the sign convention of Eq.~\eqref{eq:pairwise},
\begin{equation}
  \Var\!\left(\widehat g_i-s_{ij}\widehat g_j\right)
  =\sum_\alpha\frac{1}{n_\alpha}\Var\!\left(F_{i\alpha}-s_{ij}F_{j\alpha}\right),
\end{equation}
so the covariance that helps enters through co-measurement inside a shared context, not through the two generators sharing Pauli products. That premise was tested on the finalist pair of each system, inside the parent contexts, by comparing the cost of resolving the difference alone against the cost of resolving both generators to the same radius.

\begin{table}[!htbp]
\centering
\scriptsize
\caption{Separating the two finalists inside the shared parent contexts: resolving only their difference against resolving each of them, both to radius gap/2. All sums are over parent contexts, so these are the $S^{\text{par}}$ of Sec.~\ref{sec:breadth-depth}. The two finalist gradients have the \emph{same} sign in all four cases ($s_{ij}=+1$), so the relevant observable is $C_i-C_j$ and no sign correction applies; the one pair that is opposite-sign under the orbital phases used here, H$_4$ side $2.0$ HF, is not among them.}
\label{tab:difference}
\begin{tabular}{lrrrrr}
\toprule
Case & $S^{\text{par}}_i$ & $S^{\text{par}}_j$ & $S^{\text{par}}_{i-j}$ & Resolve both & Resolve difference \\
\midrule
H$_4$, side $1.0$, HF & 1.921 & 1.731 & 2.455 & 554,857 & 499,910 \\
LiH, $R{=}3.0$, HF & 2.170 & 2.180 & 3.623 & 48,712 & 67,577 \\
H$_2$O eq., $R{=}0.9572$, HF & 7.263 & 12.173 & 17.803 & 1,002,360 & 1,581,034 \\
H$_2$O str., $R{=}1.75$, HF & 6.239 & 7.794 & 13.005 & 10,636,302 & 18,049,318 \\
\bottomrule
\end{tabular}
\end{table}

Resolving the difference is cheaper only on H$_4$, by $1.11$; on LiH and both H$_2$O rows it is more expensive, by up to $1.7$. Writing $S$ for $S^{\text{par}}$: resolving two generators separately costs $S_i^2+S_j^2$ in units of $\epsilon^{-2}$ against $S_{i-j}^2$ for the difference, so the difference wins only when $S_{i-j}<\sqrt{S_i^2+S_j^2}$, which for comparable arms needs roughly $30\%$ cancellation. The measured cancellation $1-S_{i-j}/(S_i+S_j)$ is $33\%$ on H$_4$, $17\%$ on LiH and $7\%$ on H$_2$O stretch. It does not track the finalist overlaps of Table~\ref{tab:overlap} --- H$_4$ has zero overlap and the largest cancellation, LiH $35\%$ overlap and half the cancellation --- consistent with the premise that the useful covariance comes from co-measurement rather than shared support, but leaving the size of the cancellation unexplained.

This does not refute Part II, whose other mechanisms --- coefficient splitting over an overlapping cover, empirically estimated covariances, adaptive choice of the next context --- are untouched. It does mean the pairwise-difference mechanism should not be assumed to pay, and that cancellation must be measured on the intended benchmark first. A second caution: that mechanism acts on the depth term, only $12\%$ of M3's cost on LiH and $39\%$ on H$_2$O stretch, so its effect is bounded even where it pays.

Part III~\cite{PartIII} changes the ADAPT trajectory so that the commutators first become symmetry-filtered and tapered. When Part III is active, the Part I object $C_i$ is replaced by
\begin{equation}
  C_{i,\mathrm{tap}}^{\parallel,k}
  =\mathrm{Taper}_{\mathcal R_k}\left(\Pi_{\mathcal R_k}([H,G_i])\right),
\end{equation}
and then M1, M2, or M3 is applied to these reduced operators. Part I tables do not include this reduction.

\section{Remaining work in Part I}
\label{sec:remaining}

\begin{enumerate}[leftmargin=2em]
  \item \textbf{Adopt actual precision in the M3 accounting}, or report both
  columns. The planning bound is conservative by up to a factor of two on the
  larger systems (Sec.~\ref{sec:spillover}), and the attribution, the M2/M3
  comparison and the staged regrouping of Sec.~\ref{sec:hybrid} were all computed
  on it; the last has not been repeated on actual precision.
  \item \textbf{Allocate incrementally.} Solving each breakpoint independently of
  the shots already committed is what lets M3 exceed M1 (Sec.~\ref{sec:m3-m1}).
  Topping up from committed shots repairs the counterexample found; whether it
  guarantees $M_{\mathrm{M3}}\le M_{\mathrm{M1}}$ in general is open.
  \item \textbf{Extend the routine sweep} to the finite-shot and correlated studies,
  which are still run under one grouping routine.
  \item \textbf{Price the excluded cost components.} Distinct settings differ by
  $3.5$ to $6.1$ times between the method families (Table~\ref{tab:settings}), and
  no circuit depth is counted, so the M2-versus-M3 comparison is incomplete even
  where the shot counts are resolved.
  \item \textbf{Shot-level sampling.} The surrogate samples estimators; neither its
  Gaussian marginals nor its joint draw has been compared with multinomial sampling
  of measurement outcomes.
  \item \textbf{Establish whether the CISD state is a fair proxy} for a
  mid-trajectory ADAPT state, which needs trajectory data.
\end{enumerate}

\section{What is tested, and what is not}
\label{sec:tests}

\textbf{Covered.} 41 unit tests. Against dense linear algebra, independently of
the Pauli-string evaluator: every commutator expansion against $HG_i-G_iH$ as a
matrix, every gradient against a central finite difference of the energy, and every
fragment variance and context covariance against dense fragment matrices, all on
H$_4$ side $1.0$ HF. Conventions and closed forms: basis ordering and the
big-endian convention; the qubit Hamiltonian against the PySCF RHF and FCI
energies on every case, as a hard gate; the allocation against
Eq.~\eqref{eq:single-gradient-shot} for a single gradient; the fragment
decomposition summing back to the gradient operator. Groupings: every parent
context and every individual grouping is a partition of its support into fully
commuting sets. Invariances and pins: method totals and pairwise variances
unchanged under negating any subset of generators; both square-H$_4$ RHF energies;
the context covariances positive semidefinite, within Cauchy--Schwarz, with exact
diagonals and non-negative pairwise variances under either sign convention
(Sec.~\ref{sec:correlation}); the counterexample to $M_{\mathrm{M3}}\le
M_{\mathrm{M1}}$ of Sec.~\ref{sec:m3-m1}; and the actual-precision M3 of
Sec.~\ref{sec:spillover} as the limit of the noiseless trial.
\texttt{verify\_part1.py} adds 68 checks on the published results: each reported
total is the sum of the file that explains it, the gate passed, the allocation
meets its confidence target, $M_{\mathrm{M3}}\le M_{\mathrm{M1}}$ on every case,
all methods agree on the selected generator, repeated runs agree exactly, the
ratios are invariant to $\delta$, and a geometric schedule converges onto the exact
limit. All pass.

\textbf{Not covered}, in rough order of value: shot-level sampling of measurement
outcomes, against which neither the Gaussian marginals (assumption (i) of
Sec.~\ref{sec:assumptions}) nor the joint draw has been compared; the dense checks
on any system beyond H$_4$; KKT residuals for the allocation on systematic
multi-constraint instances; and regression pinning of every table entry against a
commit hash.

\textbf{Square H$_4$ has more than one RHF solution, and the build pins one.}
From eight perturbed starting guesses at side $1.0$~\si{\angstrom}, seven converge
to $-1.761075054$~Ha and one to $-1.694889592$~Ha; at side
$2.0$~\si{\angstrom} all eight agree on the energy to $4\times10^{-13}$~Ha but
their densities differ by up to $1.0$ in max-norm --- degenerate symmetry
partners, which relabel which generator carries which gradient. PySCF's default
guess is not a well-defined target here: with the same PySCF version, one
numerical environment converges at side $1.0$ to the lower solution and another to
the higher, which is internally unstable and whose gradients differ by up to
$0.25$. The build therefore follows any internal instability until the stability
analysis passes. This reaches $-1.761075054$~Ha in every environment tried,
reproduces the reported gradients to $10^{-15}$ and the method totals exactly, and
changes nothing on the other seven geometries, all of which are stable at the
default guess. A test pins both H$_4$ energies.

\emph{Within} one solution the statistics are phase-invariant: a sign flip changes
$g_i$ and the matching row and column of the covariance together, leaving $|g_i|$,
every fragment variance and $s_{ij}\Cov(\widehat g_i,\widehat g_j)$ --- hence
$R_{ij}$ of Eq.~\eqref{eq:Rij} --- unchanged, and a test confirms that negating an
arbitrary subset of generators leaves every method total and every pairwise
variance of Eq.~\eqref{eq:pairwise} unchanged. \emph{Across} solutions nothing is
guaranteed, not even generator labelling, which is why signed gradients and raw
covariances from different runs must not be compared.

The implementation should be regarded as validated at the level these tests reach
and no further.

\section{Reproduction}
\label{sec:reproduction}

\begin{verbatim}
python -m pip install -r requirements.txt
python scripts/run_part1_methods.py --all --out runs
python scripts/run_finite_shot.py --case H4_square_eq_side1p0_HF \
    --trials 200 --out runs
python scripts/experiment_q7_decomposition.py --case LiH_R3p0_HF
python scripts/experiment_spillover.py --out runs
python scripts/experiment_grouping_routines.py --case H2O_stretch_HF \
    --out runs --resume
python scripts/experiment_m3ext.py --case H4_square_eq_side1p0_HF
python scripts/experiment_calibration.py --case LiH_R3p0_HF
python scripts/experiment_m3_vs_m1_adversarial.py --compressed
python -m pytest tests -q && python scripts/verify_part1.py --runs runs
\end{verbatim}

Case identifiers accepted by \texttt{-{}-case} are
\texttt{LiH\_R3p0\_HF}, \texttt{H4\_square\_eq\_side1p0\_\{HF,CISD\}},
\texttt{H4\_square\_stretch\_side2p0\_\{HF,CISD\}},
\texttt{H2O\_\{eq,stretch\}\_\{HF,CISD\}}.

Each case writes a per-method summary together with the FC groups, fragment
standard deviations and BAI histories behind it, and records the interpreter and
package versions with every run. \texttt{verify\_part1.py} checks the published
results rather than the implementation, as listed in Sec.~\ref{sec:tests}; all 68
checks pass alongside 41 unit tests.

\textbf{Determinism and seeds.} The planning bounds contain no randomness and
re-run exactly; \texttt{verify\_part1.py} asserts this. The H$_2$O sweep takes
several minutes per grouping, and \texttt{-{}-resume} saves each grouping as it
completes, so an interrupted sweep loses at most the one in progress. The finite-shot tables use
\texttt{numpy} default-RNG streams at seed $0$, except Table~\ref{tab:calibration},
whose per-entry seeds are recorded in the run summaries. The routine sweep uses
\texttt{random} orderings at seeds $0$ to $3$.

\textbf{Numerical settings.} PySCF defaults are used for the SCF convergence
threshold and initial guess, except that square H$_4$ needs the phase convention
noted in Sec.~\ref{sec:tests} to make covariances reproducible. The CISD state is
the lowest eigenpair of the Hamiltonian restricted to the $\le2$-excitation space,
taken with a dense symmetric eigensolver below $2^{12}$ basis states and with
\texttt{scipy.sparse.linalg.eigsh} at tolerance $10^{-12}$ above it. The
convex-solver cross-check of Sec.~\ref{sec:proxy} used SLSQP on randomly generated
instances with up to eight arms and twenty contexts, agreeing to better than one
part in $10^{6}$.

\textbf{Environment.} Code and run outputs live at
\url{https://github.com/Zephrous5747/ADAPT_Sampling}. Pin Python 3.10.12,
PySCF 2.14.0, OpenFermion 1.8.1, NumPy 2.2.6 and SciPy 1.15.3;
\texttt{requirements.txt} states minimum versions only and does not by itself
recreate that environment. On one desktop core the wall clock is dominated by
building the gradient problem: about $16$~s for H$_4$, $37$~s for LiH, two to
three minutes for H$_2$O; each 200-trial run then takes $0.1$ to $40$~s per
method.

\begin{thebibliography}{9}
\setlength{\itemsep}{1pt}

\bibitem{Grimsley2019ADAPT}
H. R. Grimsley, S. E. Economou, E. Barnes, and N. J. Mayhall,
``An adaptive variational algorithm for exact molecular simulations on a quantum computer,''
\emph{Nature Communications} \textbf{10}, 3007 (2019).
\url{https://doi.org/10.1038/s41467-019-10988-2}

\bibitem{Shkolnikov2023Avoiding}
V. O. Shkolnikov, N. J. Mayhall, S. E. Economou, and E. Barnes,
``Avoiding symmetry roadblocks and minimizing the measurement overhead of adaptive variational quantum eigensolvers,''
\emph{Quantum} \textbf{7}, 1040 (2023).
\url{https://doi.org/10.22331/q-2023-06-12-1040}

\bibitem{Anastasiou2023ReallyMeasure}
P. G. Anastasiou, N. J. Mayhall, E. Barnes, and S. E. Economou,
``How to really measure operator gradients in ADAPT-VQE,''
\emph{arXiv:2306.03227} (2023).
\url{https://arxiv.org/abs/2306.03227}

\bibitem{Huang2025QuantumGambling}
R. Huang and A. F. Izmaylov,
``Quantum Gambling: Best-Arm Strategies for Generator Selection in Adaptive Variational Algorithms,''
\emph{arXiv:2509.14917} (2025); \emph{J. Chem. Theory Comput.} (2026).
\url{https://arxiv.org/abs/2509.14917}

\bibitem{Rossi2026RSe}
E. Rossi, E. R. Kjellgren, A. F. Izmaylov, S. P. A. Sauer, K. M. Ziems, and S. Coriani,
``Resource-efficient energy-based operator selection in fermionic ADAPT-VQE via exact Hamiltonian transformation,''
\emph{arXiv:2606.04786} (2026).
\url{https://arxiv.org/abs/2606.04786}

\bibitem{PartII}
L. Wang,
``Part II: Variance-Aware Universal Operators for Grouped Best-Arm Identification in ADAPT-VQE Gradient Measurements,''
internal project note (2026); source in \texttt{reports/} of the repository cited in Sec.~\ref{sec:reproduction}.

\bibitem{PartIII}
L. Wang,
``Part III: Symmetry-Staged ADAPT-VQE with Residual Stabilizer Tapering,''
internal project note (2026); source in \texttt{reports/} of the repository cited in Sec.~\ref{sec:reproduction}.

\end{thebibliography}

\end{document}
